\documentclass[twoside, openright, fontsize=10pt, headings=big]{scrbook}

% Load your existing preamble
% ============================================================
%  preamble.tex  —  Textbook Note-Taking Preamble & Macros
%  Include with: \input{preamble}
%  
%  Designed for use with scrbook document class
% ============================================================

% ------------------------------------------------------------
%  ENCODING & FONTS
% ------------------------------------------------------------
%\usepackage[T1]{fontenc}
%\usepackage[utf8]{inputenc}
%\usepackage{lmodern}
\usepackage{microtype}

% ------------------------------------------------------------
%  PAGE GEOMETRY
% ------------------------------------------------------------
\usepackage[
  a4paper,
  top=2.5cm, bottom=2.5cm,
  inner=2.8cm, outer=3.5cm,   % textbook-style mirrored margins
  marginparwidth=2.2cm,
  marginparsep=4mm
]{geometry}

% ------------------------------------------------------------
%  COLOUR PALETTE  (redefine here to retheme the whole file)
% ------------------------------------------------------------
\usepackage[dvipsnames,svgnames,x11names]{xcolor}

\definecolor{AccentMain}{HTML}{1A6B8A}   % deep teal  — headings, rules
\definecolor{AccentWarm}{HTML}{C0392B}   % crimson    — warnings / important
\definecolor{AccentSoft}{HTML}{27AE60}   % green      — examples / tips
\definecolor{BoxBg}     {HTML}{EAF4FB}   % pale blue  — theorem background
\definecolor{WarnBg}    {HTML}{FEF9E7}   % pale amber — caution boxes
\definecolor{CodeBg}    {HTML}{F4F4F4}   % light grey — code background
\definecolor{MarginClr} {HTML}{7F8C8D}   % muted grey — margin notes

% ------------------------------------------------------------
%  TikZ  (loaded early — needed by tcolorbox skins & titlepage)
% ------------------------------------------------------------
\usepackage{tikz}

% ------------------------------------------------------------
%  ICONS  (must be loaded BEFORE tcolorbox box definitions)
% ------------------------------------------------------------
\IfFileExists{fontawesome5.sty}{%
  \usepackage{fontawesome5}%
}{%
  \newcommand{\faExclamationTriangle}{!}%
  \newcommand{\faBookOpen}{}%
  \newcommand{\faLightbulb}{}%
}

% ------------------------------------------------------------
%  MATH PACKAGES
% ------------------------------------------------------------
\usepackage{amsmath, amssymb, amsthm, mathtools, bm}

% ------------------------------------------------------------
%  TCOLORBOX
% ------------------------------------------------------------
\usepackage{tcolorbox}
\tcbuselibrary{most}

\newcounter{tcbtheorem}[section]
\renewcommand{\thetcbtheorem}{\thesection.\arabic{tcbtheorem}}

%  Theorem
\newtcbtheorem[use counter=tcbtheorem, number within=section]
  {theorem}{Theorem}{
    enhanced, breakable,
    colback=BoxBg, colframe=AccentMain,
    fonttitle=\bfseries\color{white},
    attach boxed title to top left={yshift=-2mm, xshift=4mm},
    boxed title style={colback=AccentMain, rounded corners},
    separator sign={ --},
  }{thm}

%  Definition
\newtcbtheorem[use counter=tcbtheorem, number within=section]
  {definition}{Definition}{
    enhanced, breakable,
    colback=BoxBg!60, colframe=AccentMain!70,
    fonttitle=\bfseries\color{AccentMain},
    attach boxed title to top left={yshift=-2mm, xshift=4mm},
    boxed title style={colback=AccentMain!15, colframe=AccentMain!70},
    separator sign={ --},
  }{def}

%  Lemma
\newtcbtheorem[use counter=tcbtheorem, number within=section]
  {lemma}{Lemma}{
    enhanced, breakable,
    colback=BoxBg!40, colframe=AccentMain!50,
    fonttitle=\bfseries, separator sign={ --},
  }{lem}

%  Corollary
\newtcbtheorem[use counter=tcbtheorem, number within=section]
  {corollary}{Corollary}{
    enhanced, breakable,
    colback=BoxBg!40, colframe=AccentMain!50,
    fonttitle=\bfseries, separator sign={ --},
  }{cor}

%  Proposition
\newtcbtheorem[use counter=tcbtheorem, number within=section]
  {proposition}{Proposition}{
    enhanced, breakable,
    colback=BoxBg!40, colframe=AccentMain!50,
    fonttitle=\bfseries, separator sign={ --},
  }{prop}

%  Example
\newtcbtheorem[use counter=tcbtheorem, number within=section]
  {example}{Example}{
    enhanced, breakable,
    colback=AccentSoft!8, colframe=AccentSoft!70,
    fonttitle=\bfseries\color{AccentSoft!80!black},
    attach boxed title to top left={yshift=-2mm, xshift=4mm},
    boxed title style={colback=AccentSoft!20, colframe=AccentSoft!70},
    separator sign={ --},
  }{ex}

%  Important / Warning
\newtcolorbox{important}[1][]{
  enhanced, breakable,
  colback=WarnBg, colframe=AccentWarm,
  title={\faExclamationTriangle~Important},
  fonttitle=\bfseries\color{AccentWarm},
  attach boxed title to top left={yshift=-2mm, xshift=4mm},
  boxed title style={colback=WarnBg, colframe=AccentWarm},
  #1
}

%  Remark
\newtcolorbox{remark}[1][]{
  enhanced, breakable,
  colback=gray!7, colframe=gray!50,
  title={Remark},
  fonttitle=\bfseries\color{gray!60!black},
  #1
}

%  Key Idea
\newtcolorbox{keyidea}[1][Key Idea]{
  enhanced, breakable,
  colback=AccentMain!8, colframe=AccentMain,
  title={#1},
  fonttitle=\bfseries\color{white},
  attach boxed title to top left={yshift=-2mm, xshift=4mm},
  boxed title style={colback=AccentMain},
}

%  Intuition
\newtcolorbox{intuition}[1][Intuition]{
  enhanced, breakable,
  colback=AccentSoft!10, colframe=AccentSoft,
  title={\faLightbulb~#1},
  fonttitle=\bfseries\color{AccentSoft!80!black},
}

%  Proof style
\renewenvironment{proof}[1][\proofname]{%
  \noindent\textit{#1.}\enspace\ignorespaces
}{%
  \hfill$\blacksquare$\par\medskip
}

% ------------------------------------------------------------
%  MARGIN NOTES
% ------------------------------------------------------------
\usepackage{marginnote}
\renewcommand*{\marginfont}{\footnotesize\color{MarginClr}\itshape}
\newcommand{\mnote}[1]{\marginnote{\raggedright#1}}

% ------------------------------------------------------------
%  HEADERS & FOOTERS 
% ------------------------------------------------------------
\usepackage{titlesec}

\titleformat{\chapter}[display]
  {\color{AccentMain}\Large\bfseries}
  {\chaptertitlename\ \thechapter}{12pt}
  {\Huge\color{black}}
  [\vspace{2pt}{\color{AccentMain}\titlerule[1.5pt]}]

\titleformat{\section}
  {\color{AccentMain}\large\bfseries}
  {\thesection}{0.8em}{}
  [\vspace{1pt}{\color{AccentMain!50}\titlerule[0.6pt]}]

\titleformat{\subsection}
  {\color{AccentMain!80!black}\normalsize\bfseries}
  {\thesubsection}{0.7em}{}

\titleformat{\subsubsection}
  {\color{AccentMain!60!black}\normalsize\itshape\bfseries}
  {\thesubsubsection}{0.7em}{}

% ------------------------------------------------------------
%  HYPERREF & BOOKMARKS  (load late; cleveref must come after)
% ------------------------------------------------------------
\usepackage[
  colorlinks=true,
  linkcolor=AccentMain,
  citecolor=AccentMain!70,
  urlcolor=AccentSoft!80!black
]{hyperref}
\usepackage{bookmark}

% ------------------------------------------------------------
%  LISTS
% ------------------------------------------------------------
\usepackage{enumitem}
\setlist[itemize,1]{label=\textcolor{AccentMain}{\textbullet}, leftmargin=1.5em}
\setlist[itemize,2]{label=\textcolor{AccentMain!60}{\textendash}, leftmargin=1.5em}
\setlist[enumerate,1]{label=\textcolor{AccentMain}{\arabic*.}, leftmargin=1.8em}

% ------------------------------------------------------------
%  FIGURES & TABLES
% ------------------------------------------------------------
\usepackage{graphicx, float, caption, subcaption, booktabs, array, multirow}
\captionsetup{
  labelfont={bf, color=AccentMain},
  font=small,
  skip=4pt
}

% ------------------------------------------------------------
%  CODE LISTINGS
% ------------------------------------------------------------
\usepackage{listings}
\lstset{
  backgroundcolor=\color{CodeBg},
  basicstyle=\ttfamily\small,
  breaklines=true,
  frame=single,
  framerule=0pt,
  xleftmargin=8pt, xrightmargin=8pt,
  keywordstyle=\color{AccentMain}\bfseries,
  stringstyle=\color{AccentWarm},
  commentstyle=\color{AccentSoft!80!black}\itshape,
  numberstyle=\tiny\color{MarginClr},
  numbers=left, numbersep=6pt,
  showstringspaces=false,
  captionpos=b
}

% ------------------------------------------------------------
%  MISCELLANEOUS PACKAGES
% ------------------------------------------------------------
\usepackage{afterpage}     % \afterpage{\nopagecolor} — required by titlepage
\usepackage{lipsum}        % dummy text — remove in final doc
\usepackage{siunitx}       % \SI{9.8}{\metre\per\second\squared}
\usepackage{physics}       % \dd, \dv, \pdv, \grad, \curl, \div, etc.
\usepackage{cancel}        % \cancel{x}
\usepackage{cleveref}      % \cref — must be loaded after hyperref

% ------------------------------------------------------------
%  MATH MACROS
% ------------------------------------------------------------

%  ---- Number sets ----
\newcommand{\N}{\mathbb{N}}
\newcommand{\Z}{\mathbb{Z}}
\newcommand{\Q}{\mathbb{Q}}
\newcommand{\R}{\mathbb{R}}
\newcommand{\C}{\mathbb{C}}
\newcommand{\F}{\mathbb{F}}

%  ---- Set-builder notation ----
\newcommand{\st}{\,:\,}                    % such that
\newcommand{\given}{\,\vert\,}             % conditional bar
\newcommand{\setof}[2]{\left\{#1\st#2\right\}}

%  ---- Math operators ----
\DeclareMathOperator{\sgn}{sgn}
\DeclareMathOperator{\lcm}{lcm}
\DeclareMathOperator{\Span}{span}
\providecommand{\rank}{\relax}
\renewcommand{\rank}{\operatorname{rank}}
\DeclareMathOperator{\nullity}{nullity}
\providecommand{\trace}{\relax}
\renewcommand{\trace}{\operatorname{tr}}
\DeclareMathOperator{\diag}{diag}
\DeclareMathOperator{\proj}{proj}
\DeclareMathOperator{\Hom}{Hom}
\DeclareMathOperator{\End}{End}
\DeclareMathOperator{\Aut}{Aut}
\DeclareMathOperator{\im}{im}
\DeclareMathOperator{\coker}{coker}

%  ---- Linear algebra ----
\newcommand{\mat}[1]{\mathbf{#1}}          % bold matrix/vector
\newcommand{\T}{^{\mathsf{T}}}             % transpose
\newcommand{\inv}{^{-1}}                   % inverse
% physics defines \norm{} and \abs{} — override with our delimiters
\renewcommand{\norm}[1]{\left\lVert#1\right\rVert}
\renewcommand{\abs}[1]{\left\lvert#1\right\rvert}
\newcommand{\inner}[2]{\left\langle#1,\,#2\right\rangle}
\newcommand{\ceil}[1]{\left\lceil#1\right\rceil}
\newcommand{\floor}[1]{\left\lfloor#1\right\rfloor}

%  ---- Calculus / analysis ----
\newcommand{\eps}{\varepsilon}
\newcommand{\Del}{\Delta}
% physics redefines \lim, \limsup, \liminf — restore ours
\renewcommand{\lim}{\operatorname*{lim}}
\renewcommand{\limsup}{\operatorname*{lim\,sup}}
\renewcommand{\liminf}{\operatorname*{lim\,inf}}
\newcommand{\dint}[2]{\int_{#1}^{#2}}

%  ---- Probability & statistics ----
\newcommand{\E}[1]{\mathbb{E}\!\left[#1\right]}
\newcommand{\Var}[1]{\mathrm{Var}\!\left(#1\right)}
\newcommand{\Cov}[2]{\mathrm{Cov}\!\left(#1,#2\right)}
\newcommand{\Prob}[1]{\mathbb{P}\!\left(#1\right)}
\newcommand{\iid}{\overset{\text{iid}}{\sim}}
\newcommand{\dist}[1]{\operatorname{#1}}

%  ---- Logic / notation ----
\newcommand{\then}{\Rightarrow}
\renewcommand{\iff}{\Leftrightarrow}
\newcommand{\conj}[1]{\overline{#1}}
\newcommand{\comp}[1]{{#1}^{\mathsf{c}}}

%  ---- Misc ----
\newcommand{\RNum}[1]{\uppercase\expandafter{\romannumeral #1\relax}}

% ------------------------------------------------------------
%  KOMA-SCRIPT COMPATIBILITY NOTES:
%  - The titlesec package may cause warnings with scrbook.
%  - For full KOMA compatibility, consider replacing titlesec
%    with KOMA's native commands (setkomafont, addtokomafont, 
%    chapterprefix=false, etc.)
%  - Headers/footers are configured in the main .tex file
%    using scrlayer-scrpage (fancyhdr has been removed)
% ------------------------------------------------------------

% ============================================================
%  END OF PREAMBLE
% ============================================================
\usetikzlibrary{cd}
% KOMA header configuration
\usepackage{scrlayer-scrpage}
\clearpairofpagestyles
\automark[chapter]{chapter}
\ohead{\pagemark}
\ihead{\headmark}
\cehead{\MakeUppercase{\TextbookTitle}}
\renewcommand{\chapterpagestyle}{empty}
\newcommand{\addunumberedsubsec}[1]{%
    \subsection*{#1}%
    \addcontentsline{toc}{subsection}{#1}%
}

% Define title page fields
\def\YourName{Deepak Jassal}
\def\TextbookTitle{Structure of Groups and Rings}
\def\TextbookAuthor{Dr. Stanley Xiao}
\def\CourseName{MATH420}
\def\NoteSemester{Winter}
\def\NoteYear{2026}

% Make equations share the exact same counter as theorems
\makeatletter
\let\c@equation\c@tcbtheorem
\renewcommand{\theequation}{\thetcbtheorem}
\makeatother


\begin{document}
%No page spaces between chapters
\let\cleardoublepage\clearpage
% Title page
% ============================================================
%  titlepage.tex  —  Detailed Textbook Notes Title Page
% ============================================================

\begin{titlepage}
  \centering
  
  \vspace*{1.5cm}
  
  % Top decoration
  {\Large \scshape \textbf{Personal Study Notes} \par}
  
  \vspace{0.3cm}
  \rule{0.3\textwidth}{0.4pt}
  \vspace{1cm}
  
  % Course info
  {\Large \CourseName \par}
  
  \vspace{0.5cm}
  
  % Main title
  {\Large \textsc{Based on} \par}
  {\Huge \bfseries \TextbookTitle \par}
  
  \vspace{0.5cm}
  
  % Textbook author
  {\Large \itshape by \TextbookAuthor \par}
  
  \vspace{1.5cm}
  
  \rule{0.5\textwidth}{0.4pt}
  
  \vspace{1cm}
  
  % Your info
  {\Large \textbf{Notes by:} \YourName \par}
  
  \vspace{0.3cm}
  
  % Date info
  {\large \NoteSemester\ \NoteYear \par}
  
  \vfill
  
  % Footer
  \rule{0.8\textwidth}{0.4pt}
  \vspace{0.2cm}
  
  {\small \textit{These notes are for personal use only.}}
  
\end{titlepage}

\newpage
\thispagestyle{empty}
\mbox{}
\newpage

% ============================================================
%  END OF TITLE PAGE
% ============================================================

\frontmatter
\tableofcontents
\mainmatter

\chapter{Lecture 1}
\addunumberedsubsec{Rings}
Rings are a particular kind of algebraic structure.\\
What is the point in considering ``algebraic structures''?\\
Algebraic structures are abstractions that clarify some essential feature of a problem, and provide a language to describe issues and how they may be resolved.\\
\\
Rings ``come from'' groups.\\
Groups are in some sense the simplest kind of abstract algebraic structures.\\
Groups are however, already \underline{very} hard structures to understand.\\
However groups simply cannot capture many objects we come about, which fundamentally may have more than one operation.\\
From a group $G=(G,\ast)$ we can always form a ring called the \underline{group-ring}.\\
\begin{definition}{Rings}{}
    A set $\mathcal{R}$ together with two binary operations $``+''$ and $``\times''$ such that
    \begin{enumerate}
        \item $(\mathcal{R},+)$ is an Abelian group;
        \item $\times$ is associative;
        \item $a\times(b+c)=a\times b + a\times c,$ $\forall a,b,c\in \mathcal{R}$\\;
              $(b+c)\times a=b\times a + c\times a$ $\forall a,b,c\in \mathcal{R}$\\;
    \end{enumerate}
    forms a ring $\mathcal{R}=(\mathcal{R},+,\times)$.\\
    A ring with a multiplicative identity $1_\mathcal{R}$ such that $1_\mathcal{R}\times a=a\times1_\mathcal{R}=a$ is called a \underline{unital} ring.
\end{definition}
\begin{example}{}{}
    \begin{enumerate}
        \item $(\Z,+,\times)$ is a unital ring,
        \item $(2\Z,+,\times)$ is a ring, but it is not a unital ring.
    \end{enumerate}
    Both are examples of \underline{commutative} rings, because $\times$ commutes in both cases.
\end{example}
\begin{definition}{}{}
    $
        M_{2\times2}(\Z)=\left\{ \begin{bmatrix}
            1&b\\c&d
        \end{bmatrix}: a,b,c,d\in\Z \right\}
    $
    with respect to matrix adidtion and multiplication is a non-commutative ring.
\end{definition}
The jury is out on whether non-commutative rings are important. Heavily context dependant.
\chapter{Lecture 2}
Last time: Defined a ring  and introduced same examples of commutative and non-commutative rings.\\
Two most important examples are polynomial rings and matrix rings. Polynomial rings are commutative and matrix rings are non-commutative rings. These two types of rings make up a large portion of rings which are studied.
\begin{example}{}{}
    Let $S\subseteq\C$ be a (possibly finte) set, and $f:\C\mapsto\C$ be a meromorphic function. Then $f$ is regular on $\C\setminus S$ if f has no poles outside of $S$.\\
    Let $R(S)$ be the set of functions regular on $\C\setminus S$. $R(S)$ forms a ring with respect to function addition and multiplication.
\end{example}
\addunumberedsubsec{One Point Compactification}
Place the unit sphere into the complex plane. Then, if we use the point (0,0,1) we can uniquely map each point in $\C$ to a point on the unit sphere. All but one point on the unit sphere is mapped to, the missing point is (0,0,1). If we append $\{\infty\}$ to $\C$ we can make this function a bijection between the unit sphere and $\C\cup\{\infty\}$.\\
However interesting this might be it makes it so the only holomorphic functions are the constand functions.
\begin{example}{}{}
    The following is an exmaple of a quotient ring,
    \[
        \R[x,y]=\{p(x,y):p\text{ is a polynomial with real coefficients}\}.
    \]
    One can obtain the quotient ring
    \[
        \underbrace{\R[x,y]\setminus(x^2-y)}_{\text{Corresponds to } y=x^2}
    \]
\end{example}
\begin{example}{}{}
    For any $n\in\N$ the set $T_n(\C)$ of upper triangular (respectively lower triangular) $n\times n$ matrices forms a ring with respect to matrix addition and matrix multiplication.
\end{example}
\begin{remark}
    Matrices have a group structure under rotation, however reflections break this structure.
\end{remark}
Symmetries are usually described by certain matrix groups, called \textit{Lie groups}. Associated to Lie groups are \textit{Lie algebras}. These are \underline{matrix rings}.
\[
    \text{Lie group}\underset{\text{exponential map}}{\overset{\frac{d}{dx}}{\xrightleftharpoons{\hspace{2cm}}}}\text{Lie algebras}
\]
\begin{example}{}{}
    Number rings i.e.,
    \[
        \Z[\sqrt{2}]=\{a+b\sqrt{2}:a,b\in\Z\}\subseteq\R.
    \]
    The addition and multiplication is inherited from $\R$.
    To see that this is a ring we need only to check that it closed under the adidtion and multiplication inherited from $\R$, the other properties are given from $\R$.\\
    Let $a_1+b_1\sqrt{2},a_2+b_2\sqrt{2}\in\Z[\sqrt{2}]$. Then,
    \[
        a_1+b_1\sqrt{2}+a_2+b_2\sqrt{2}=\underbrace{a_1+a_2}_{\in\Z}+\underbrace{(b_1+b_2)}_{\in\Z}\sqrt{2}.
    \]
    \begin{align*}
        (a_1+b_1\sqrt{2})(a_2+b_2\sqrt{2})&=a_1a_2+a_1b_2\sqrt{2}+a_2b_1\sqrt{2}+\sqrt{2}^2b_1b_2\\
        &=\underbrace{a_1a_2+2b_1b_2}_{\in\Z}+\underbrace{(a_1b_2+a_2b_1)}_{\in\Z}\sqrt{2}.
    \end{align*}
    Not all number rings have unique factorization (ex. $\Z[\sqrt{-5}]$).\\
\end{example}
\begin{remark}
    We refer to the ``dimension'' of a ring as the rank of a ring.
\end{remark}
\begin{center}
    Integers $\leftrightarrow$ quadratic rings\\
    Binary cubic forms acted on by $GL_2(\Z)$ $\leftrightarrow$ cubic rings 
\end{center}
\begin{remark}
    Binary cubic forms:$ax^3+bx^2+cxy^2+dy^3,a,b,c,d\in\Z.$  
\end{remark}


\chapter{Lecture 3}
\textit{Recall.} Qoutients in ring theory were briefly mentioned.\\
The original motivation for this phenomenon came from modular arithmetic.\\
Consider the (OG ring) $(\Z,+,\times)$, We can consider an arithmetic induced by a positive integer $m$, through remainders when divied by m.\\
\begin{example}{}{}
    If $m=3$, every integer can be written as $3x$, $3x+1$ or $3x+2$.
    \[
      (3x+1)(3y+2)=\underbrace{9x+6x+3y}_{3z}+2.  
    \]
\end{example}
The modern way to define modular arithmetic modulo $m$:
\begin{enumerate}
    \item Introduce an equivalence relation $\sim_m$ on $\Z$:\\
    $a\sim_m b$ \textit{if and only if} $a-b$ is divisible by $m$.
    \item Introduce addition and multiplication to the equivalence classes with respect to $\sim_m$
    \[
        [a]_m+[b]_m:=[a+b]_m,
    \]
    \[
        [a]_m[b]_m:=[a\times b]_m.
    \]
    \item Check this is well defined. That is if the representative of a class changes, output does not chance under addition or multiplication.
\end{enumerate}
Then we get a new ring $(\Z/m\Z,+,\times)$.\\\\
Can this be done in general?\\
If the ring is commutative we get a very nice thoery. If the ring is not commutative this becomes much more complicated; more interpretations.\\
Even for commutative rings quotients of the form $R/S$ with $S$ an arbitrary subring are not necessarily well behaved.\\
For $R/S$ to inherit the ring structure on $R$ we need that $[a]_S[b]_S=[ab]_S$ to be well defined. If there exists $s\in S$ and $r\in R$, then
\[
    [r]_S\underbrace{[s]_S}_{=[0]_S}=[rs]_S\neq0.
\]
Thus $S$ needs to satisfy a much stronger condition, that being $sr\in S$ for all $s\in S$ and $r\in R$. Such subrings are called ``ideals''.
\begin{definition}{Ideals}{}
    Let $(R,+,\times)$ be a commutative ring. A subset $S\subset R$ togeth with $+,\times$ is said to be \underline{ideal} if 
    \begin{enumerate}
        \item $(S,+)$ is an Abelian group;
        \item for all $s\in S$ and $r\in R$ we have $sr\in S$. 
    \end{enumerate}
\end{definition}
\begin{proposition}{}{}
    Let $(R,+,\times)$ be a commutative ring and $I\subseteq R$ an ideal. Then $R/I,+,\times$ is a ring.
\end{proposition}
\begin{proof}
    \textit{Sketch}.\\
    We have to show that 
    \[
        [a_1]_I[b_1]_I=[a_2]_I[b_2]_I.
    \]
    Hypothesis implies that there exists $s_1,s_2\in I$ such that $a_2=a_1+s_1$ and $b_2=b_1+s_2$. Then,
    \begin{align*}
        [a_2]_I[b_2]_I&=[a_2b_2]_I\\
        &=[(a_1+s_1)(b_1+s_2)]_I\\
        &=[a_1b_1+s_1b_1+s_2a_1+s_1s_2]_I\\
        [a_1]_I[b_1]_I&=[a_1b_1]_I\qedhere
    \end{align*}    
\end{proof}
To study rings, just like vector spaces we want to understand the \underline{maps} between rings.
\begin{definition}{Ring Homomorphisms}{}
    Let $(R,+,\times)$,$(S,\oplus,\otimes)$ be rings.\\
    A function $f:R\mapsto S$ is said to be a homomorphism If
    \begin{enumerate}
        \item $f(0_R)=0_S$;
        \item $f(a+b)=f(a)\oplus f(b)$;
        \item $f(ab)=f(a)\otimes f(b)$.
    \end{enumerate}
\end{definition}
\begin{important}
    If $R,S$ are commutative rings, and $f:R\mapsto S$ is a homomorphism, then the kernel of $f$ is an ideal.
\end{important}

\chapter{Lecture 4}

\textit{Last time:} Ring Homomorphisms
\[
    f: R\to S
\]
\[
    f(a+b)=f(a)+f(b)
\]
\[
    f(ab)=f(a)f(b)
\]
\[
\ker(f)=\{r\in R: f(r)=0\}
\]
What properties for $\ker(f)$ have?\\
In linear algebra, $\ker(f)$ is a \underline{subspace}. In group theory we know that it is a normal \underline{subgroup}.
\begin{lemma}{Ring Homomorphism Kernel's}{}
    $\ker(f)$ ($f$ is a homomorphism) is a subring of $R$.
\end{lemma}
\begin{proof}
    Suppose $a,b\in\ker(f)$.
    \begin{enumerate}
        \item $a+b\in\ker(f)$
        \[
            f(a+b)=f(a)\oplus f(b)=0_S\oplus 0_S=0_S.
        \]
        \[
            f(ab)=f(a)\otimes f(b)=\underbrace{0_S\otimes 0_S=0_S}_{\substack{\text{requires proof} \\ \text{Excersize}}}.
        \]
        \begin{align*}
            f(-a)&=f((-1)a)\\
            &=f(-1)\otimes \underbrace{f(a)}_{\in \ker(f)}\\
            f(-1)\otimes 0_S&=0_S
        \end{align*}
        \begin{align*}
            f(0_R)&=0_S\\
            f(a+(-a))&=f(a)\oplus f(-a)\\
            \Rightarrow f(-a)=-f(a)&=-0_S=0_S.\qedhere
        \end{align*}
    \end{enumerate}
\end{proof}
\begin{definition}{Subring}{}
    Let $(R,+,\times)$ be a ring. A subset $S\subseteq R$ is said to be a subring of $R$ if $(S,+,\times)$ is a ring.
\end{definition}
Looking to group theory, kernels are not just subgroups by \underline{normal} subgroups.
\begin{proposition}{}{}
    Let $R,S$ be rings and $f:R\to S$ be a homomorphism. Then $\ker(f)$ is an ideal of $R$. (Note: if $R$ is not commutative, then a (two-sided) ideal $I$ is one where $sr,rs\in I$ for all $s\in I$ and $r\in R$).    
\end{proposition}
\begin{proof}
    Already proved that it is a subring. Suffice to prove that for all $s\in I=\ker(f)$ and $r\in R$, $rs,sr\in I$.
    \[
        f(sr)=f(s)\otimes f(r)=0_S,
    \]
    \[
        f(rs)=f(r)\otimes f(s)=0_S.\qedhere
    \]    
\end{proof}
Since $\ker(f)$ is an ideal, $R\setminus\ker(f)$ is a (quotient) ring.\\
In linear alegbra, there is the \textit{rank-nullity} theorem:
\begin{theorem}{Rank-Nullity Theorem}{}
    If $V,W$ are vector spaces over the same scalar field $F$, and $L:V\to W$ a linear map, then $\mathrm{dim}(L(V))+\mathrm{dim}(\ker(L))=\mathrm{dim}(V).$
\end{theorem}
The analogy of this is the \underline{First Isomorphism Theorem}.
\begin{theorem}{First Isomorphism Theorem}{}
    Let $R,S$ be rings and $f:R\to S$ a homomorphism. Then $R\setminus\ker(f)\cong f(R)$. Here $\cong$ means isomorphic.
\end{theorem}
\begin{definition}{Isomorphic}{}{}
    A ring homomorphism is an isomorphism if it is invertible and the inverse map is also a homomorphism.
\end{definition}
\begin{lemma}{Homorphic Image of Rings}{}
    The homorphic image of a ring is a ring.
\end{lemma}
\begin{proof}[Proof for Lemma]
    Suppose $x,y\in f(R)$, by definition there exists $a,b\in R$ such that $f(a)=x$, and $f(b)=y$.
    \[
        x\oplus y=f(a)\oplus f(b)=f(a+b)\Rightarrow x\oplus y\in f(R).
    \]
    \[
        x\otimes y=f(a)\otimes f(b)=f(a\times b)\Rightarrow x\otimes y\in f(R).
    \]
    Prove there exists an inverse map as an excersize. \qedhere
\end{proof}
\begin{proof}[Proof of the First Isomorphism Theorem]
   Recall $r\setminus\ker(f)$ is a set of cosets (or representatives of equivalence classes). $[a]_{\ker(f)}:=a+\ker(f).$ We want to write down a map, say $\tau$, mapping each coset to an element in $f(R)$.\\
   Define 
   \[
        \tau:r\setminus\ker(f)\to f(R)
   \]
   by $\tau(a+\ker(f))=f(a)$.\\
   Check well-definedness:\\
   If $a-a'\in\ker(f)$, then 
   \begin{align*}
        \tau(a+\ker(f))&=\tau(a+(a-a')+\ker(f))\\
        f(a)&=\tau(a'+\ker(f))\\
        &=f(a').\qedhere
   \end{align*}
\end{proof}

\chapter{Lecture 5}
\textit{Last Time:} Defined a \underline{canonical} isomorphism (candidate) $\tau:R\setminus\ker(f)\to f(R)$.
\[
    \tau(a+\underbrace{\ker(f)}_{+I})=f(a)
\]
is well defined.\\
\[
    \tau((a+I)+(b+I))=\tau((a+b)+I)=f(a+b)=f(a)\oplus f(b)=\tau(a+I)\oplus\tau(b+I),
\]
\[
    \tau((a+I)(b+I))=\tau((ab)+I)=f(a+b)=f(a)\otimes f(b)=\tau(a+I)\otimes\tau(b+I).
\]
Thus, $\tau$ is a homomorphism.\\
Must show that $\tau$ is bijective.\\
If $y\in f(R)$, then there exists $e\in R$ such that $f(a)=y$. Hence, $\tau(a+I)=f(a)=y$. Thus, $\tau$ is surrjective.\\
If $\tau(a+I)=\tau(b+I)$, then 
\begin{align*}
    0_S&=\tau(a+I)\oplus(-\tau(b+I))\\
    &=f(a)\oplus f(b)\\
    &=f(a)\oplus f(-b)\\
    &=\tau((a-b)+I)\Rightarrow a-b\in I,
\end{align*}
thus, $a+I=b+I$. Hence, $\tau$ as a function is bijective. Therefore, $\tau^{-1}$ is well defined.\\
It remains to show that $\tau^{-1}$ is a homomorphism. Given that $y_1=f(a_1)$ and $y_2=f(a_2)$
\begin{align*}
    \tau^{-1}(y_1\oplus y_2)&=\tau^{-1}(f(a_1)\oplus f(a_2))\\
    &=\tau^{-1}(f(a_1+a_2))\\
    &=\tau^{-1}(\tau((a_1+a_2)+I))\\
    &=(a_1+a_2)+I\\
    &=\tau^{-1}(y_1)+\tau^{-1}(y_2)
\end{align*}
Similar process for multiplication. Thus, we have that $\tau^{-1}$ is a homomorphism.\\
Ideals were actually meant to be ``numbers''. In fact, ``Ideal numbers'' are a thing. This compels and \underline{arithmetic} on the set of ideals.
\begin{definition}{Arithmetic on Ideal}{}
    Let $R$ be a ring, and $I,J$ (2-sided) ideals. Then:
    \[
        I+J:=\{i+j:i\in I,j\in J\}
    \]
    \[
        I\cap J=\{m:m\in I\cap J\}
    \]
    \[
        I\cdot J:=\left\{\sum_{\ell=1}^{q}j_\ell i_\ell:j_\ell\in J, i_\ell\in I \right\}
    \]
\end{definition}
\begin{proposition}{}{}
    Let $R$ be a ring and $I,J$ ideals of $R$. Then $I+J$, $I\cap J$ and $I\cdot J$ are ideals. 
\end{proposition}
\begin{proof}
    \begin{enumerate}
        \item $I+J$ is an ideal. \\
        Suppose $a_1,a_2\in I+J$. Then there exist $i_1,j_1,i_2,j_2$ such that $a_k=i_k+j_k$ for $k=1,2$. Then
        \begin{align*}
            a_1+a_2&=(i_1+j_1)+(i_2+j_2)\\
            &=\underbrace{(i_1+i_2)}_{\in I}+\underbrace{(j_1+j_2)}_{\in J}\\
            &\in I+J.
        \end{align*}
        \begin{align*}
            a_1a_2&=(i_1+j_1)(i_2+j_2)\\
            &=i_1i_2+i_1j_2+i_2j_1+j_1j_2\\
            &\in I+J.
        \end{align*}
        Other details are an exercise.
        \item $I\cap J$ is an ideal.\\
        Suppose $a_1,a_2\in I\cap J$. Then $a_1+a_2\in I$ and $a_1+a_2\in J\Rightarrow a_1+a_2\in I\cap J$. Similarly $a_1a_2\in I$ and $a_1a_2\in J\Rightarrow a_1a_2\in I\cap J$. Other details are an exercise.
        \item $I\cdot J$ is an ideal.
        $a_1,a_2\in I\cdot J$.
        \[
            a_1=\sum_{\ell=1}^{q_1}i_\ell j_\ell,
        \]
        \[
            a_2=\sum_{\ell'=1}^{q_2}i_\ell' j_\ell',
        \]
        then
        \[
            a_1+a_2=\sum_{\ell=1}^{q_1}i_\ell j_\ell+\sum_{\ell'=1}^{q_2}i_\ell' j_\ell'.
        \]
        Define $q=\max\{q_1,q_2\}$, WLOG $q=q_1$, define $i_\ell'=j_\ell'=0$ if $q_2<\ell \leq q_1$. then
        \[
            a_1+a_2=\sum_{\ell=1}^{q}(i_\ell j_\ell+i_\ell'+j_\ell').
        \]
        \textit{Note.} $i_\ell j_\ell+i_\ell'+j_\ell'\in I\cdot J$ for $1\leq \ell\leq q$. Hence, $a_1+a_2\in I\cdot J$.
        \begin{align*}
            a_1a_2&=\left(\sum_{\ell=1}^{q_1}i_\ell j_\ell\right)\left(\sum_{\ell'=1}^{q_2}i_\ell' j_\ell'\right)\\
            &=\sum_{\ell=1}^{q_1}i_\ell j_\ell\left(\sum_{\ell=1}^{q_2}i_\ell'j_\ell'\right)\\
            &=\sum_{k=1}^{q_1}\sum_{\ell=1}^{q_2}i_k j_k i_\ell'j_\ell'.
        \end{align*}
        Other details are an exercise.
    \end{enumerate}
    This is a finite sum of summands of the form $ij$. Hence, $a_1a_2\in I\cdot J$.
\end{proof}
(If $r\in R$, $ra_2=\sum_{\ell=1}^{q}ri_\ell'j_\ell'\Rightarrow ra_2\in I\cdot J$)\\
Original Motivation: trying to solve diophantine equations (Fermat's last theorem).
\addunumberedsubsec{Fermat's Equation}
\[
    x^n+y^n=z^n,\quad x,y,z\in\Z,n\geq3
\]
\textit{Claim (FLT).} This has no solutions with $xyz\neq0$.\\
Fermat solved the case for $n=4$. He saw that you can factor the LHS as $(x+\zeta_1y)\cdots(x+\zeta_ny),$ where $\zeta_i$'s are the $n-$th roots of unity. If $x+\zeta_iy$ were all $n-$th powers, then FLT would be easy to prove. Unfortunately this is not case in general. This led to the notion of ``ideal numbers'' to try to recover unique factorization. This ``recovered'' by a theorem of Hilbert. Unfortunately, in practice it does not help.\\
This is because the arithmetic of ideals just introduced does not behave well with with relevant number theoretic properties. The multiplicative aspect of ideal arithmetic gave rise to so-called \underline{composition laws} in number rings.\\
This is realted to the \textit{class number problem}. Gauss had two conjectures:
\begin{enumerate}
    \item For imaginary quadratic fields, the class numbers go to infiity, and only a specified list of such fields have class \#1 (i.e., unique factorization)
    \item For real quadratic fields, infiitely many have unique factorization. 
\end{enumerate}
1. was proven by showing that if the Reimann hypothesis is true then it is true, and if the Reimann hypothesis is false then the statement is also true.

\chapter{Lecture 6}
\begin{theorem}{Second Isomorphism Theorem}{}
    Let $R$ be a ring and $I,J$ ideals of $R$. Then $I$ is an ideal of $I+J$, and $I\cap J$ is an ideal of $J$ then
    \[
        I+J\setminus J\cong I\setminus I\cap J
    \]
\end{theorem}
\begin{proof}[Proof of the Second Isomorphism Theorem]
    Clearly $I\subseteq I+J$, and also a subring. For $r\in I+J\Rightarrow r\in R$, so $r\cdot i\in I$ for $i,r\in I$ for all $i\in I$, since $I\subseteq R$ is an ideal in $R$. Similar argument shows $I\cap J$ is an ideal in $J$. Now define
    \[
        \tau:I+J\to J\setminus I\cap J,
    \]
    \[
        \tau(i+j)=i+(I\cap J)=[I]_{I\cap J}.
    \]
    It is possible that $i_1+j_1=i_2+j_2$ so we must check that this is well-defined.
    \[
        \underset{\in I}{i_1-i_2}=\underset{\in J}{j_2-j_1}\Rightarrow i_1-i_2,j_2-j_1\in I\cap J.
    \]
    \[
        \tau(i_1+j_2)=i_1+(I\cap J)=i_1-(i_1-i_2)+(I\cap J)=i_2+(I\cap J)=\tau(i_2+j_2).
    \]
    Next, we must show that $\tau$ is a homomorphism.
    \begin{align*}
        \tau(i_1+j_2+i_2+j_2)&=\tau(i_1+i_2+j_1+j_2)\\
        &=i_1+i_2+(I\cap J)\\
        &=i_1+(I\cap J)+i_2+(I\cap J)\\
        &=\tau(i_1+j_1)+\tau(i_2+j_2),
    \end{align*}
    \begin{align*}
        \tau((i_1+j_1)(i_2+j_2))&=\tau(i_1i_2+i_1j_2+j_1i_2+j_1j_2)\\
        &=i_1i_2+(I\cap J)\\
        &=(i_1+(I\cap J))(i_2+(I\cap J))\\
        &=\tau(i_1+j_1)\tau(i_2+j_2).
    \end{align*}
    Since $\tau(i+j)=i+I\cap J$ we have that $\tau$ is surrjective.\\
    Now we need only show that $\ker(\tau)=J$.\\
    Suppose $j\in J$. Then $\tau(0+j)=0+(I\cap J)$, so $j\in\ker(\tau)$ ($J\subseteq\ker(\tau)$).\\
    Now suppose that $j\in\ker(\tau)$. Then $\tau(i+j)=i+I\cap J=0+(I\cap J)$. Hence, $i\in (I\cap J)$, to $i\in J$. Therefore, $i+j\in J$. So $\ker(\tau)\subseteq J$. Thus, $J=\ker(\tau)$. Hense, by the first isomorphism theorem we have our desired result. 
\end{proof}
\begin{theorem}{Third Isomorphism Theorem}{}
    Let $R$ be a a ring and $I,J$ ideals of $R$ such that $I\subseteq J$. Then,
    \[
        R\setminus J\cong \frac{R\setminus I}{J\setminus I}.
    \]
\end{theorem}
\begin{proof}[Proof of the Third Isomorphism Theorem]
    Left as an excersize.
\end{proof}
The earliest\footnote{Probably the earliest} rings to be considered are close cousins of $\Z:$ (integral) domains.
\begin{definition}{Integral Domains}{}
    A ring $(D,+,\times)$ is siad to be an \textit{integral domain} if
    \begin{enumerate}
        \item $D$ is unital (has multiplicative identity $1\neq 0$);
        \item $D$ is commutative;
        \item for $a,b\in D$, if $ab=0$, then $a=0$ or $b=0$.
    \end{enumerate}
\end{definition}

\chapter{Lecture 7}
For fields, a key invariant if the \underline{characteristic}. $\mathrm{Char}(F)=$ smallest positive integer $p$ such that $\underbrace{1+\cdots+1}_{p\text{ times}}$. If such $p$ exists, $\mathrm{Char}(F)=0$. The problem arises when one considers fields with ``mixed characteristic''.\\
Fields with the same mixed characteristic play nice with each other, enabling various \underline{cohomology theories} (\'etale, de Rham, crystalline, prismatic).\\
If a field of one element does exist, it would sync up all of these difference cohomology theories ``somehow''. That's why such an object is unlikely.\\
\begin{definition}{Prime Ideals}{}
    Let $(D,+,\times)$ be a domian. We say an ideal $I\subseteq D$ is \underline{prime} if for all ideals $A,B\subseteq D$, if $A\cdot B\subseteq I$, then $A\subseteq I$ or $B\subseteq I$.
\end{definition}
\begin{definition}{Generated Ideals}{}
    Let $R$ be a ring and $S\subseteq R$. The ideal generated by $S$ denoted $\langle S\rangle$ is the smalled ideal containing $S$.
\end{definition}
\begin{example}{}{}
    Let $D=\Z$. $(k)-$ ideal generated by $\{k$\}. $p\mid ab$ means every multiple of $ab$ is also a multiple of $p$. This implies that $(ab)\subseteq(p)$, the first is multplies of $ab$ and the second is all multiples of $p$.
\end{example}
\begin{definition}{Maximal Ideals}{}
    An ideal $I\subsetneqq R$ is \underline{maximal} if whenever $J\subseteq R$ is an ideal and $I\subseteq J$, then $I=J$ or $J=R$.
\end{definition}
\begin{proposition}{Maximal Ideals in $\Z$}{}
    If $D=\Z$ then every prime ideal is maximal.
\end{proposition}
\begin{proof}
    Let $p$ be a prime. Then $(p)$ is a prime ideal. We want to show that it is maximal. Suppose $(p)\subseteq J$. We may assume that $J\neq\Z$. Then $J=(k)$ (because $\Z$ is a PID) for some $k\in \Z$. Then ``every multiple of $p$ is a multiple of $k$''. In particular, $p$ is a multiple of $k$. Hence $k=\pm p,\pm1$. If $k=\pm1$, $J=R=\Z$ which violates the assumption. Hence $k=\pm p$ so $(k)=(p)$.
\end{proof}
\begin{proposition}{}{}
Let $(D,+,\times)$ be a domain. Then every maximal ideal $I\subseteq D$ is prime.
\end{proposition}
\begin{proof}
    Suppose $I$ is maximal but not prime. Then there exists $A,B\subseteq D$ ideals such that  $AB\subseteq I$ but $A\nsubseteq I$ and $B\nsubseteq I$. There exists $a\in A\setminus I$ and $b\in B\setminus I$, but $ab\in I$.\\
    Let $J=\langle I\cup \{a\}\rangle$, $K=\{t+ar:i\in I,r\in D\}$. $I\subset J$, but $I\neq K$ (because $a\neq\in I$). Hence, $K=D$. Hence, $1\in K$. Therefore, there exists $i\in I$, $r\in D$ such that $1=i+ar$. Thus, $b1=b(i+ar)=bi+bar$, $bi\in bI\subseteq DI=I$. Then $bar\in I\Rightarrow b\in I$ and this is a contradiction.
\end{proof}
\begin{theorem}{}{unital1}
    Let $R$ be a unital, commutative ring. Let $I\subseteq R$ be an ideal. Then $1\Leftrightarrow 2$, whenever
    \begin{enumerate}
        \item $R\setminus I$ is an integral domain;
        \item $I$ is prime.
    \end{enumerate}
\end{theorem}
\begin{theorem}{}{unital2}
    Let $R$ be a unital, commutative ring. Let $I\subseteq R$ be an ideal. Then $1\leftrightarrow2$, whenever
    \begin{enumerate}
        \item $R\setminus I$ is a field
        \item $I$ is \underline{maximal}.
    \end{enumerate}
\end{theorem}
Theorem \ref{thm:unital2} to \ref{thm:unital1} is how we get that all maximal ideals in an itegral domain are prime. All field's are integral domains.

\chapter{Lecture 8}

\textit{Last Time:}
\begin{theorem}{}{}
    Let $R$ be a unital commutative ring and $I\subseteq R$ an ideal of $R$. Then $I$ is prime (proper) \textit{iff} $R/I$ is an integral domain.
\end{theorem}
\begin{proof}[Proof of the Theorem]
    Suppose $I$ is a prime ideal. Then $R/I$ is a ring. For $r,s\in R$ we have 
    \begin{align*}
        (r+I)(s+I)&:=rs+I\\
        &=sr+I\\
        (s+I)(r+I).
    \end{align*}
    So $R/I$ is also commutative. Also
    \begin{align*}
        (1+I)(r+I)&=1r+I\\
        &=r+I,
    \end{align*}
    so $1+I$ is the identity of $R/I$. Must show that $(r+I)(s+I)=0+I$. Since $rs\in I$ and $I$ is prime we have $r\in I$ or $s\in I$, as required.\\
    Now assume that $R/I$ is an integral domain. Then $R/I$ is commutative unital and has no zero-divisors. \textit{Note.} $1+I$ is a multiplicative identity in $R/I$, so it must be the only multiplicative identity since it is unique. Since additve and multiplicative identities are never equal we have
    \[
        1+I\neq 0+I.
    \]
    In particular $I\neq R$.\\
    Suppose $AB\subseteq I$ with $A,B$ ideals of $R$. Thus, for all $a\in A$ and $b\in B$ we have $ab\in I$. Suppose for the sake of a contradiction, that $A\not\subset I$ and $B\not\subset I$. Then we can choose $a\in A\setminus I$ and $b\in B\setminus I$, such that $ab\in I$. Now
    \begin{align*}
        (a+I)(b+I)&=ab+I\\
        &=0+I.
    \end{align*}
    So if $(a+I)$ and $(b+I)$ are non-zero in $R/I$ the must be zero divisors, this is a contradiction. Thus, $I$ is a prime ideal.
\end{proof}
\begin{theorem}{}{}
    Let $R$ be a unital, commutative ring, $I\subseteq R$ be a proper ideal. Then $R/I$ is a field \textit{iff} $I$ is maximal. 
\end{theorem}
\begin{proof}[Proof of the Theorem]
    Suppose $I$ is maximal. Then $I$ is prime. Hence $R/I$ is an integral domain. Remains to prove non-zero elements in $R/I$ are invertible.\\
    Suppose $r\in R$ such that $r+I\neq 0+I$ (i.e., $r\neq I$). Put $N=\{sr+t:s\in R, t\in I\}$. Let $a\in R$. Then $a(sr+t)=(as)r+at\in N$. Hence $N$ is an ideal of $R$. But $1\cdot r+0\in N$, but $r\not\in I$. Hence $I\subsetneqq N$. Therefore $N=R$. Hence there exists $s,t$ such that $sr+t=1$. For this chouse fo $s,t$ we have 
    \begin{align*}
        (s+I)(t+I)&=sr+I\\
        &=(sr+t)+I\\
        &=1+I.
    \end{align*}
    Hence, $(s+I)$ is the inverse of $(t+I)$ in $R/I$.\\
    Now suppose that $I$ is not maximal. If $I$ is not prime then $R/I$ is not an integral domain, hence not a field. Hence we may assume that $I$ is prime. Since $I$ is not maximal, there exists an ideal $J$ such that $J\subsetneqq R$ and $I\subsetneqq J$. In particular $i\not\in J$. Let $a\in J/I$. If $ab+I=1+I$ for $b\in R$ then $ab-i\in I$. Hence $ab-1=i$ for some $i\in I$. Note $ab\in J$. Then $i\in J$, since $I\subset J$. Therefore, $1=ab-i\in J$ which is a contradiction.
\end{proof}
\begin{example}{}{}
    Consider $R=\Z[\sqrt{-5}]=\{a+b\sqrt{-5}:a,b\in\Z\}$. This is a subring of $\C$. $2,3$ are \underline{irreducible} in the sense that they cannot be written as the product of ``smaller'' elements.\\
    i,e.,
    \[
        2=\underbrace{(a+b\sqrt{-5})}_{\in R}\underbrace{(a-b\sqrt{-5})}_{\in R}\underbrace{r}_{\in\R}
    \] 
    \[
        r=(a^2+5b^2).
    \]
    Then
    \[
        \underbrace{r(a-b\sqrt{-5})}_{\in R}\Rightarrow ra,rb\in\Z.
    \]
    Hence $r\in\Q$.\\
    $r=\frac{p}{q},q>0,(p,q)=1$. then $q\mid a$ and $q\mid b$; replace $a$ with $qa$ and $b$ with $qb$ (abuse of notation). Hence,
    \begin{align*}
        2&=(qa+qb\sqrt{-5})(pa-pb\sqrt{-5})\\
        &=pq(a+b\sqrt{-5})(a-b\sqrt{-5})\\
        &-pq(a^2+5b^2).
    \end{align*}
    \[
        p=2,q=1=a^2+5b^2\Rightarrow b=0,a+b\sqrt{-5}=\pm1.
    \]
    \[
        q=2,p=1\text{ same as above.}
    \]
    \[
        p=q=1,a^2+5b^2=2
    \]
    this has no solutions. Hence, 2 is irreducible. Is it prime?\\
    2,3 are both irreducible. Are they prime? Well,
    \[
        6=2\cdot 3=(1+\sqrt{-5})(1-\sqrt{-5}).
    \]
    Consider the principal ideals
    \[
        (2),\quad (3),\quad (1+\sqrt{-5}),\quad (1-\sqrt{-5}).
    \]
    None of these ideals are prime.  
\end{example}

\chapter{Lecture 9}
\begin{example}{}{}
    $R=\Z[\sqrt{-5}]$\\
    \textit{Claim:} $(2),(3),(1+\sqrt{-5}),(1-\sqrt{-5})$ are not prime (as ideals). We have that $(2)$ is prime \textit{iff} $R/(2)$ is an integral domain.\\
    $R/(2)$ is a set of equivalence classes.\\
    i.e., $[1]_{(2)}=1+(2),[0]_{(2)},[\sqrt{-5}]_{(2)},[1+\sqrt{-5}]_{(2)}$ are these all the equivalence classes?\\
    Consider $[a+b\sqrt{-5}]_{(2)}$ with $a,b\in\Z$. If $a\equiv b\equiv 0\mod2$ then $[a+b\sqrt{-5}]_{(2)}=[0]_{(2)}$. If $a$ is odd, and $b$ is even we get $[a+b\sqrt{-5}]_{(2)}=[1]_{(2)}$. If $a$ is even and $b$ is odd then $[a+b\sqrt{-5}]_{(2)}=[\sqrt{-5}]_{(2)}$. If $a$ is odd and $b$ is odd then $[a+b\sqrt{-5}]_{(2)}=[1+\sqrt{-5}]_{(2)}$. Thus we have $R/(2)=\{[0]_{(2)},[1]_{(2)},[b\sqrt{-5}]_{(2)},[1+\sqrt{-5}]_{(2)}\}$. 
    \[
        [1+\sqrt{-5}]_{(2)}[1+\sqrt{-5}]_{(2)}=[0]_{(2)}.
    \]
    Since there are zero divisors we have that $R/(2)$ is not an integral domain. Thus, $(2)$ is not prime.\\
    $R/(3)=\{[0],[1],[2],[1+\sqrt{-5}],[1+2\sqrt{-5}],[2+\sqrt{-5}],[2+2\sqrt{-5}],[\sqrt{-5}],[2\sqrt{-5}]\}$
    \[
        [1+2\sqrt{-5}][1+\sqrt{-5}]=[1+3\sqrt{-5}-10]=[0].
    \]
    $p=11$, $R/(p)$. Let $a_1,b_1\in\{0,1,\dots,10\}$. Want $[a_1+b_1\sqrt{-5}]_{(11)}$ is invertible, provided that $a_1,b_1$ are not both zero.\\
    \textit{Consider}: $[a_2+b_2\sqrt{-5}]_{(11)}$ then
    \[
       [a_1+b_1\sqrt{-5}]_{(11)}[a_2+b_2\sqrt{-5}]_{(11)}=[a_1a_2-5b_1b_2+\sqrt{-5}(a_1b_2+a_2b_1)]\tag{$\ast$}.
    \]
    Easy if one or the other is zero. Choose $(a_2,b_2)=k(a_1-b_1)$. Then ($\ast$) becomes $[k(a_1^5)+5b_1^2]_{(11)}$. Then $a_1^5+5b_1^2$ is divisible by 11 \textit{iff} $a_1,b_1$ are both divisible by 11.
\end{example}
Since $(2),(3)$ are \underline{not} prime ideals so they should be ``divisible'' by some other ideal. In particular, not prime implies not maximal. So there is some ideal $J_2,J_3$ such that $(2)\subsetneqq J_2$ and $(3)\subsetneqq J_3$. Extend (2) with $1+\sqrt{-5}$ to get $\langle 2,1+\sqrt{-5} \rangle$.\\
$R/\langle2,1+\sqrt{-5} \rangle=\{[0],[1]\}=\mathbb{F}_2$. Thus, the ideal is maximal and therefore prime. 
\[
    J_2\cdot J_2=\{\sum_{j=1}^{2}a_jb_j:a_j,b_j\in J_2\}.
\]
$a=2x+y(1+\sqrt{-5})$, $b=2u+v(1+\sqrt{-5})$. Then
\[
    a\cdots b=0+(2),
\]
Thus
\[
    J_2\cdot J_2=(2).
\]
$(3)=\langle3,1+\sqrt{-5} \rangle\langle3-\sqrt{-5} \rangle$ and neither of these ideals are the whole ring.\\
\textit{Summary:} $(2)=J_2^2$ and $(3)=J_3\overline{J_3}$, ($J_3=\langle3,1+\sqrt{-5} \rangle$, $\overline{J_3}=\langle3-\sqrt{-5} \rangle$).\\
$(5)$ is also a square. (Excersize).\\
(11) is a prime ideal.\\
$(41)=(6+\sqrt{-5})(6-\sqrt{-5})$.

\chapter{Lecture 10}
\subsection*{Why Rings?}
We want to do computations with rings. So we want to be able to run the original algorithm: the \underline{Euclidean algorithm}.
\begin{definition}{Euclidean Norm}{}
    Let $D$ be an integral domain. A Euclidean norm (valuation) is a function $\mu:D\to\Z_{\geq0}$ such that 
    \begin{enumerate}[label=(\alph*)]
        \item for all $a,b\in D$ with $b\neq0$, there exists $q,r\in D$ such that
        \[
            a=qb+r
        \]
        and either $\mu(r)=0$ or $\mu(r)<\mu(b)$
        \item and for all $a,b\in D\setminus\{0\}$, $\mu(a)\leq\mu(ab)$.
    \end{enumerate}
\end{definition}
Part (a) of the definition induces a division algorithm on $D$.\\
There is a division algorithm, and Euclidean algorithm on $\R[x]$. The valuation is given by $\deg(f)$.
\begin{theorem}{Eulclidean Domain}{}
    A domain $D$ norm $\mu$ is a \underline{Euclidean domain}.
\end{theorem}
\begin{example}{}{}
    $\Z$, $\R[x],$ $\Z[\sqrt{-1}]$ are Eulclidean domains, and $\Z[\sqrt{-5}]$ is not a Eulclidean domain.    
\end{example}
\begin{important}
    Impotant: Eulciean domain $\subset$ principal ideal domain (PID) $\subset$ unique factorization domains (UFD).
\end{important}
Towards Euclidean algorithm.
\begin{definition}{Divisors}{}
    Let $R$ be a commutative ring. Suppose $a,b\in R$, $b\neq0$. We say $b$ divides $a$ or $b\vert a$ if there exists $q\in R$ such that $a=qb$. Otherwise we say $b\nmid a$.
\end{definition}
\begin{theorem}{GCD}{}
    Let $a,b\in R$ we say $g\in R$ is a greatest common divisor of $a,b$ if $g\vert a$ and $g\vert b$, and for all $c\in R$ such that $c\vert a$ and $c\vert b$ we have $c\vert g$.
\end{theorem}
\begin{theorem}{The Extended Euclidean Algorithm}{}
    Let $D$ be a Euclidean domain and $a,b\in D$. Then $a,b$ has a gcd in $A$, say $g$. Further there exist $x,y\in D$ such that 
    \[
        g=ax+by.
    \]
\end{theorem}
\begin{proof}[Proof of the Theorem]
    Put $I=\{ax+by:x,y\in D\}=(a)+(b)$ which is an ideal. Using the fact that $D$ being a Eulclidean domain implies that $D$ is a PID (will be proved later). There exists $d\in D$ such that $I=(d)$. Then there exists $p,q\in D$ such that $d=ap+bq$. Note that $a\in I$ and $b\in I$. Hence there exists $k_1,k_2\in I$ such that $a=dk_1$ and $b=dk_2$. Hence $d$ is a common divisor of $a,b$. Any common divisor of $a,b$ must also divide $app+bq$, so $d$ is a gcd. Suppose that $g$ is another gcd of $a,b$ then we have $g\vert d$ and $g\vert d$. Hence $d=gm$ and $g=dn$ for some $m,n\in D$. $d=(dn)m=d(mn)$
\end{proof}

\chapter{Lecture 11}

\begin{theorem}{Euclidean Domains $\Leftrightarrow$ PID}{}
    Euclidean domains are principal ideal domains.
\end{theorem}
\begin{proof}[Proof of the Theorem]
    Let $T\subset D$ be an ideal. If $I=\{0\}$, then $I=(0)$. Otherwise, $I$ contains a non-zero element. Let $S_I=\{\mu(i):i\in I\}\subseteq\Z_{\leq0}$. By the well-ordering principal $S_I$ has a least element. Hence there exists $k\in I$ such that $\mu(k)=\min\{S_I\}$.
    
    \begin{claim}{}{clm:generator}
        $I=(k)$.
    \end{claim} 
    
    \begin{proof}[Proof of the Claim]
        Suppose $\ell\in I$, $\ell\neq0$. By the division algorithm (part a) there exists $q,r\in D$ such that 
        \[
            \ell=qk+r,\quad \mu(r)<\mu(k),\quad r=0.
        \]
        \[
            \ell\in I\Rightarrow \ell-pk=r\in I.
        \]
        By definition $\mu(k)=\min\{S_I\}$ so $r=0$. Therefore $\ell=pk\in(k)$, so $I\subseteq(k)$. The other direction is trivial.
    \end{proof}
\end{proof}
\begin{definition}{Unique Factorization Domains}{}
    Let $D$ be an integral domain.
    \begin{enumerate}[label=(\alph*)]
        \item An element $q\in D$ is a unit if $q$ has a multiplicative inverse.
        \item An element $q\in D$ is said to be prime if $(q)$ is a prime ideal.
        \item An element $q\in D$ is said to be irreducible if $q=ab$, $a,b\in D$ then one of $a,b$ is a unit.
        \item If $q\in D$ and $u\in D$ is a unit, then $qu$ is an associate of $q$. 
    \end{enumerate}
\end{definition}
\begin{proposition}{}{}
    Let $D$ be a domain. If $q\in D$ is prime, then it is irreducible
\end{proposition}
\begin{proof}
    Let $q\in D$ be prime. Suppose for the sake of a contradiction, that $q$ is reducible. Then there exists non-units $a,b\in D$ such that 
    \[
        q=ab.
    \]
    Since $q$ is a prime, $(q)$ is a prime ideal. By definition $(a)\subseteq (q)$ or $(b)\subseteq(q)$. WLOG suppose $(a)\subseteq(q)$. Then there exists $c\in D$ such that $a=qc$. Then we have $q=ab=(qc)b=q(bc)\Rightarrow q=(1-bc)=0$. Since $D$ is a domain, it has no zero divisors. Since 0 is not prime we assume that $q\neq0$, so $1-bc=0$ but $b$ is not a unit.
\end{proof}
\begin{theorem}{}{}
    Let $D$ be a PID. Then $q\in D$ is (a) prime \textit{iff} $q$ is (b) irreducible \textit{iff} $(q)$ is (c) maximal.
\end{theorem}
\begin{proof}[Proof of the Theorem]
    (a)$\Rightarrow$(b) Always true.\\
    (b)$\Rightarrow$(a) Suppose $q$ is irreducible. Suppose $A,B\in D$ are ideals such that $AB\subseteq (q)$. Since $D$ is a PID, $A=(a)$, and $B=(b)$ for some $a\in A$ and $b\in B$. $\Rightarrow$ $ab\in(q)$, so there exists $s\in D$ such that $ab=qs$. Put $J=\langle a,q \rangle$. Since $D$ is a PID $J=\langle j\rangle$ for some $j\in D$. Hence there exists $x,y\in D$ sick that $a=jx$, and $q=jy$. Since $q$ is irreducible either $j$ or $y$ is a unit. 
    \begin{enumerate}
        \item If $j$ is a unit. Then $J=D$. Hence there exists $v,w\in D$ such that $1=av+qw\Rightarrow b=(ab)v+qbw$, $ab=qs$, so $b=(qs)v+(ab)w=q(sv+bw)\in(q)\Rightarrow (b)\subset(q)$.
        \item $y$ is a unit. Then $q,j$ are associates. This implies that $(q)=(j)=J=(a,q)$. Then we have $(a)\subseteq\langle a,q\rangle\subseteq(q)$. Hence, $q$ is prime. 
    \end{enumerate}
    (b)$\Rightarrow$(c) Suppose that $q$ is irreducible. Let $J$ be an ideal, such that $(q)\subseteq J$. Since $D$ is a PID we have $J=(j)$ for some $j\in D$. So there exists $s\in D$ such that $q=js$. Then since $q$ is irreducible $j$ or $s$ is a unit. If $j$ is a unit then $J=D$. Otherwise if $s$ is a unit then $q,j$ are associates and $(q)=(j)$. By definition $q$ is maximal.\\
    (c)$\Rightarrow$(a) We always have maximal implies prime.
\end{proof}
\begin{corollary}{}{}
    Let $\mathbb{F}$ be a field. Then and $f(x)\in\mathbb{F}[x]$ is irreducible \textit{iff} $(f(x))$ is maximal
\end{corollary}
\begin{definition}{UFD}{}
    Let $D$ be a domain. We say that $D$ is a UFD if for all $q\in D$ which is not a unit, we have:
    \begin{enumerate}[label=(\alph*)]
        \item $q=q_1\dots q_r$ where each $q_i$ $1\leq i\leq r$ are irreducible;
        \item if $q=q_1\dots q_r$ and $q=p_1\dots p_\ell$, then there exists a bijection $\sigma:\{1,\dots,r\}\to\{1,\dots,\ell\}$ such that 
        \[
            q_i=p_{\sigma(i)}U_{\sigma(i)}.
        \]
        Where $U_{\sigma(i)}$ are units.  
    \end{enumerate}
\end{definition}
\begin{definition}{Noetherian Rings}{}
    A ring $R$ is \underline{Noetherian}, or $R$ satisifies the ASC (Ascending Chain Condition), if whenever 
    \[
        J_1\subseteq\cdots\subseteq J_n\subseteq\cdots,\quad J_i\text{'s are ideals}.
    \]
    Then there exists $N\geq 1$ such that for all $m\geq N$ we have $J_m=J_N$. 
\end{definition}

\chapter{Lecture 12}
\addunumberedsubsec{Ascending Chain Condition}
If $J_1\subseteq J_2\subseteq\cdots\subseteq J_n\subseteq\cdots$ are ideals in $R$, then there exists $N\geq1$ such that $J_n=J_N$ for all $n\geq N$.
\begin{example}{}{}
    $\R[x]$ (is a PID) satisfies the ascending chain condition. $J_i=(f_i)$, ($f_i\in\R[x]$). $(f_1)\subseteq(f_2)\subseteq\cdots\subseteq(f_n)\subseteq\cdots$. $\{f_jg:g\in\R[x]\}=(f_j)\subseteq(f_{j+1})=\{f_{j+1}g:g\in\R[x]\}$. In particular, $f_j$ is a multiple of $f_{j+1}$ ($f_{j+1}\mid f_j$). So ACC follows from well ordering principle on $\Z_{\geq0}$.
\end{example}
\begin{example}{}{}
    $\R[x_1,\dots,x_n,\dots]$. $(x_1)\subseteq(x_1,x_2)\subseteq(x_1,x_2,x_3)\subseteq\cdots$.
\end{example}
\begin{theorem}{Every PID is Noetherian}{}
    If $R$ is a PID, then it is Noetherian.
\end{theorem}
\begin{proof}[Proof of the Theorem]
    Let $J_1\subseteq J_2\subseteq\cdots\subseteq J_n\subseteq\cdots$ be an AC. Put $J=\bigcup_{n=1}^\infty$. We claim that $J$ is an ideal of $R$. Suppose $a,b\in J$ and $r\in R$. Since $a,b\in J$ there exists indices $n_1,n_2$ such that $a\in J_{n_1}$, and $b\in J_{n_2}$. Without loss of generallity assume that $n_1\leq n_2$. Then $J_{n_1}\subseteq J_{n_2}$, so $a\in J_{n_2}$, Hence $a+b\in J_{n_2}$ since $J_{n_2}$ is an ideal. Thus $a+b\in J$. Now $ar\in J_{n_2}$ hence $ar\in J$. Since $R$ is a PID, there exists $\delta\in R$ such that $J=(\delta)$. In particular $\delta\in J$. Hence there exists $N\geq1$ such that $\delta\in J_N$. It follows that $J=(\delta)\subseteq J_N$. But $J_N\in J$, hence $J=J_N$.
\end{proof}
\begin{corollary}
    Let D be a PID. Suppose $d\in D$, $d\neq0$. Then there exists $P_1,\dots,P_k\in D$, $P_j$ irreducible for $1\leq j\leq k$, such that $d=P_1\cdots P_k$    
\end{corollary}
\begin{lemma}{}{}
    Let D be a PID. Suppose $d\in D$, $d\neq0$. Then $d$ has an irreducible divisor, if $d$ is not a unit.
\end{lemma}
\begin{proof}
    Suppose for the sake of a contradiction that every divisor of $d$ is reducible. Since $d\mid d$ we have that $d$ is reducible, hence $d=a_1b_1$ with $a_1,b_1$ non units. $a_1$ is reducible so $a_1=a_2b_2$, $a_2,b_2$ non units. In general $a_i=a_{i+1}b_{i+1}$. Put $J_k=(a_k)$. Then $J_k\subseteq J_{k+1}$ for all $k\geq 1$. Applying ACC, choose $N$ such that $J_k=J_N$ for all $k\geq N$. Then $J_{N+1}=J_N$. Then $(a_{N+1})\subseteq (a_N)$. Thus there exists $c\in D$ such that  $a_{N+1}=a_Nc=(a_{N+1}b_{N+1})c$. This implies that $1=b_{N+1}c$, but $b_{N+1}$ is a unit, this is a contradiction.
\end{proof}
The second step (for the corollary) proceeds by induction, and uses ACC to obtain a terminal state.\\
Factorization in polynomial rings

\chapter{Lecture 13}
Factorization in a polynomial ring.\\
Last time we proved that is $R=F$ is a field, then $F[x]$ is a Euclidean domain. In particular, there is a division algorithm on $F[x]$, which can be used to implement the EA. \\
Factoring polynomials over $R[x]$ with $R$ a ring if of significant interest. Over $\Z[x],$ one can write down an algorithm as follows:
\[
    f(x)=a_nx^n+\cdots+a_1x+a_0
\]
we factor $a_0$. For any divisor $d$ of $a_0$ and degree $r\leq n$, we look for solutions of the form $g(x)h(x)=f(x)$ where $\deg(g)=r$, and $g(x)=\cdots+g_1x+d$. This depends on prime factorization over $\Z$ which is super-polynomial time. Instead of a deterministic algorithm, how about a \underline{probablistic} algorithm? 
\begin{example}{}{}
    Is $x^2+1$ irreducible? $\Z\to\Z/3\Z$ induces a homomorphism from $\Z[x]$ to $(\Z/3\Z[x])$. $x^2+1\mod 3$ factors mod 3? Possible linear factors:
    \[
        x,2x,x+1,x+2,2x+1,2x+2.
    \]
    If $f(x)\mod p$ is irreducible over $\Z/p\Z$ for any prime $p$ then $f(x)$ is irreducible over $\Z$.
\end{example}
\begin{example}{}{}
    $x^2+x+41$, mod 2 $x^2+x+1$, mod 3 $x^2+x+2$, doesn't work for 3. Then it doesn't work till 41.     
\end{example}
Pick your favourite integers $a_1,\dots,a_n$. Let $g(n)=(x-a_1)\cdots(x-a_n)$. One can show that there exists an \underline{irreducible} $f\in\Z[x]$, such that $f(x)\equiv g(x)\mod p$ all small primes.\\

Q: If $f(x)$ is \underline{reducible} for all $p$, does it follow that $f$ is reducible over $\Z$?\\
Ans: No. $x^4+1$ is reducible for every prime $p$, but not reducible over $\Z$.\\
Assuming the Grand/Generalized Riemann Hypothesis (GRH), there is a positive number $C$ such that is $f$ is irreducible over $\Z$, there is a prime $p\leq C(\log H)^2$ such that $f(x)\mod p$ splits completely. Here 
\[
    H=H(f)=\max_{0\leq k\leq n}|f_k|
\]

\chapter{Lecture 14}
\addunumberedsubsec{Modules}
They are generalizations of vector spaces.
\begin{definition}{Modules}{}
    Let $R$ be a ring. An Abelian group $(M,\oplus)$ is said to be a left-$R$ module with respect to an operation
    \[
        \dot:R\times M\to M
    \]
    if given $r\in R, x,y\in M$
    \begin{enumerate}
        \item $r\dot(x\oplus y)=r\dot x\oplus r\dot y$,$(r+ y)\dot x=r\dot x\oplus y\dot x$;
        \item $r\dot(x\dot y)=(r\dot x)\dot y$;
        \item $1\dot x=x$.
    \end{enumerate}
\end{definition}
All Abelian groups are $\Z-$modules.
\begin{example}{}{}
    $\Z^3$ is a $\Z-$module.\\
\end{example}
Any subgroup of $(\Z^n,+)$ is a $\Z-$module, such a subgroup is denoted by $\Lambda$, say (lattice). The rank of $\Lambda$ is the max number of linearly independant vectors in $\Lambda$.\\
$\rank(\Z^3)=3$.\\
Q: Given $\Lambda\subset\Z^n$, compute a basis of a $\Lambda$. i.e., $B\subset\Lambda$ such that $\spn_{\Z}(B)=\Lambda$, $|B|=\rank(\Lambda)$.\\
Shortest basis problem for $\Lambda\subset\Z^n$ is NP-hard.\\
If $\Lambda=\spn_Z(B),B=\{\vec{b_1},\dots,\vec{b_n}\}\subset\Z^n$, then a theorem of Minkowski assers that there exists a possitive number $f(a)$ such that the shortest vector $\vec{v}\in\Lambda$ has length at most $f(n)(\det\Lambda)^{\frac{1}{n}}$, $\det\Lambda=|\det(B)|$.
\begin{definition}{Rank of a Ring}{}
    A commutative ring is said to have an \underline{order} rank of $r$ if $(R,+)$ is isomorphic to $\Z^r$ as a $\Z-$module, $R$ is unital.
\end{definition}
Q: Given an order $R$ of a finite rank, decide if $x^2+1=0$ has a solution in $R$.\\
(Lenstra) This can be decided with an algorithm in P. But to decide $x^2+x+1=0$ is NP hard. 
\begin{example}{}{}
    Elliptic Cruves
    \[
        E:y^2=x^3+Ax+B,\quad A,B\in\Z,\quad 4A^3+27B^2\neq0.
    \]
    \[
        E(\Q)=\{(x,y)\in\Q^2:y^2=x^3+Ax+B\}.
    \]
    Mordell (1922): $E(\Q)\cong G\oplus\Z^r$, for some $r\geq0$, $G$ a finite Abelian group, as groups. $G$ is the torsion part and $\Z^r$ is the free part. The $r$ in Mordell's theorem is the (algebraic) rank of $E$. 
\end{example}
Tate module.\\
Fix a prime $p$ (usually $p\nmid16(4A^3-27B^2)$). Over $\C$, $E(\C)$ has a $p-$torsion point. $E(\C)$ is a group, ``torsion'' in a group means elements of finite order.\\
Let $E[P]$ denote the set of $p-$torsion points in $E(\C)$.\\
We can play the same game for $p^n,$ $n\geq1$. Get $E[p^n]$. Turns out $E[P^n]$ are \underline{algebraic} for all $p^n$.
\addunumberedsubsec{Algebras}
\begin{example}{}{}
    Set of polynomials with real coefficients is an $\R-$algebra. $\R[x]$ clearly has the same structure of a $\R-$vector space is also a ring. 
\end{example}
\begin{definition}{$A$-Algebra}{}
    Let $A$ be a ring. An $A$-algebra is a ring $B$ together with a homomorphism $i_B:A\to B$.
\end{definition}
\begin{definition}{}
    Let $B,C$ be $A-$algebras. A homomorphism $\phi:B\to C$ is a map such that $\phi(i_B(a))=i(a)$, for all $a\in A$, where $\phi$ is a ring homomorphism.
\end{definition}

\chapter{Lecture 15}
\begin{definition}{}{}
    A ring homomorphism $i_B:A\to B$ is said to be \underline{finite-type}, and $B$ a \underline{finitely generated} $A$-algebra, if $B$ is generated by a finite set of elements as an $A$-algebra. 
\end{definition}
Compare this definition to the definition of finite-dimensional vector space $V$ over a field $F$. That being
\[
    V=\spn_F\{v_1,\dots,v_k\}
\]
for some $v_1,\dots,v_k\in V$.\\
Interpreting this for an $A$-algebra, we want:\\
there exists $b_1,\dots,b_k\in B$ such that for all $b\in B$, there exists a polynomial $P_b\in A[x_1,\dots,x_k]$ such that $b=i_B(P_b)(b_1,\dots,b_k)$.\\
Unlike vector spaces, the size of a minimal generating set (i.e., a generating fet $S$ such that $S\setminus\{s\}$ is not a generating set for all $s\in S$) of a general $A$-algebra $B$ is \underline{not} the same. So ``dimension'' is not well defined.\\
A natural family of algebras (they can be thought of as general $A$-algebras, of $F$-algebras, for $F=\Z,\Q,\R,\C$,etc.) are so-called \underline{rings of polynomial} invariants/covariants. If we want to study polynomials, we can consider how htey behave when acted upon by a group.\\
Consdier the group
\[
    GL_2(R)=\left\{\begin{bmatrix}
        a_{11}&a_{12}\\a_{21}&a_{22}
    \end{bmatrix}:
    a_{11}a_{22}-a_{12}a_{21}\neq0\right\}
\]
with $R$ a unital, commutative ring. Let 
\[
    P_2(R)=\{f_2x^2+f_1xy+f_0y^2:f_2,f_1,f_0\in R\}.
\]
$GL_2(R)$ acts on $P_2(R)$ via the \underline{substitution action}:
\[
    g=
    \begin{bmatrix}
        a_{11}&a_{12}\\a_{21}&a_{22}
    \end{bmatrix}\to f\in P_2(R),
\]
by 
\[
    g\cdot f=f(a_{11}x+a_{21}y,a_{12}x+a_{22}y)=f_2(a_{11}x+a_{21}y)^2+f_1(a_{11}x+a_{21}y)(a_{12}x+a_{22}y)+f_0(a_{12}x+a_{22}y)^2.
\]
Restrict to the subgroup $G=GL_2^{\pm1}(R)$ ($\det=\pm1$).\\
Q: does there exists a \underline{non-zero} polynomial, $Q(f)=Q(f_2,f_1,f_0)$ such that for all $g\in G$, $Q(g\cdot f)=Q(f)$?\\
Yes, $Q$ can be taken to be the \underline{discriminant}, $\Delta(f)=f_1^2-4f_2f_0$.\\
Let's try this for $G=GL_2(\Z)=GL_2^{\pm1}(\Z)$. 
\[
    G=\left\langle
    \begin{bmatrix}
        1&1\\0&1
    \end{bmatrix},
    \begin{bmatrix}
        0&1\\-1&0
    \end{bmatrix},
    \begin{bmatrix}
        0&1\\1&0
    \end{bmatrix}
    \right\rangle=\langle E_1,S,T\rangle.
\] 
\[
    E_1\cdot f=f_2x^2+f_1(x)(x+y)+f_0(x+y)^2=(f_2+f_1+f_0)x^2+(f_1+2f_0)xy+f_0y^2,
\]
\[
    \Delta(E_1\cdot f)=(f_1+2f_0)^2-4(f_2+f_1+f_0)f_0=f_1^2-4f_2f_0=\Delta(f).
\]
\[
    S\cdot f=f_2(-y^2)+f_1(-y)(x)+f_0x^2=f_0x^2-f_1xy+f_2y^2
\]
\[
    \Delta(S\cdot f)=(-f_1)^2-4f_0f_2=f_1^2-4f_2f_0=\Delta(f).
\]
Inv$_2(\Z)$ si the set of polynomials $Q(F)$ such that $Q(g\cdot f)=Q(f)$ for all $g\in G$. We just showed that $\langle\Delta(f)\rangle\subseteq$Inv$_2(\Z)$ (they are actually equal).
\[
    P_3(R)=\{f_3x^3+f_2x^2y+f+1xy^2+f_0y^3\}.
\]
$G$ acts on $P_3(R)$ by substitution. 
\[
    \Delta_3(f)=f_2^2f_1^2-4f_3f_1^3-4f_0f_2^3-27f_3^2f_0^2+18f_3f_2f_1f_0.
\]
$\Delta_3(g\cdot f)=\Delta_3(f)$ for all $g\in G$. Inv$_3(\Z)=\langle\Delta_3(f)\rangle$.\\
However for quartics we have, Inv$_4(\Z)=\langle I(f),J(f)\rangle$, with $I(f)=12f_4f_0-3f_3f_1+f_2^2$ and $J(f)=72f_4f_2f_0+9f_3f_2f_1-27f_4f_1^2-27f_0f_3^2-2f_2^3$.\\
No longer true for quintics. 
\[
    \text{Inv}_5=\langle H_6,h_8,H_{10},H_{15}\rangle/ n
\]
\[
    H^2_{15}=P(H_6,H_8,H_{10})
\]

\chapter{Lecture 16}

\addunumberedsubsec{Radical of an Ideal}
\begin{definition}{}{}
    Let $R$ be a ring, $U\subseteq R$ an ideal. The \underline{radical} of the $I$, denotes rad$(I)$, is givne by
    \[
        \text{rad}(I)=\{r\in R:r^k\in I,\text{ for some } k\geq1\}.
    \]
\end{definition}
\begin{example}{}{}
    If $R=\Z$, $I\subseteq\Z$ an ideal, what is rad$(I)$.
\end{example}
\begin{example}{}{}
    rad(3)=(3), and rad(9)=rad(3)=(3)\\
    Suppose $a\in$rad$(3)$, then there exists $k\geq1$, such that $a^k\in(3)$. Hence, $a^k=3\ell$ for some $\ell\in\Z$. Hence $3\mid a,$ for $a\in(3)$.
\end{example}
\begin{important}{}{}
    For $a\in\Z$, $a=p_1^{s_1}\cdots p_k^{s_k}$, then rad$(a)=(p_1\cdots p_k)$. Proof is left as an excersize.    
\end{important}

\begin{definition}{}{}
    For $a\in\Z$ the radical of just the integer rad$a=p_1\cdots p_k$.
\end{definition}
Q: Suppose that $a,b,c\in\Z$ such that $a+b=c$ and $(a,b,c)=1$. How large can $\max\{|a|,|b|,|c|\}$ be compared to rad$abd$?
\begin{proposition}{}{}
    For all $\varepsilon>0$, $a,b,c$ as given above, we have
    \[
        \max\{|a|,|b|,|c|\}\ll_\varepsilon\text{rad}(abc)^{1+\epsilon}.
    \]
\end{proposition}
Q: Let $R$ be a ring. What is rad(0)?\\
A: Nilpotent elements of $R$.
\begin{proposition}
    Let $R$ be a commutative ring. Then for all $I\subseteq R$ ideals, rad$I$ is an ideal. 
\end{proposition}
\begin{proof}
    Suppose $a,b\in$rad$I$. Then there exists $k,\ell\geq1$ such that $a^k,b^\ell\in I$. Then 
    \[
        (a+b)^{k+\ell}-\sum_{j=0}^{k+\ell}\binom{k+\ell}{j}a^jb^{k+\ell-j}\in I
    \]
    this is the binomial theorem in a commutative ring.\\
    Suppose $r\in R$ and $(ra)^k=r^la^k\in I$. Hence, rad$(I)$ is an ideal.
\end{proof}
\begin{lemma}{Radical of a Radical}{}
    Radicals are idempotent, that is
    \[
        \text{radrad}(I)=\text{rad}I.
    \]
\end{lemma}
\begin{proof}
    Excersize.
\end{proof}
\begin{proposition}{}{}
    Let $I\subseteq R$ be an ideal. Then 
    \[
        \text{rad}I=\bigcap_{I\subseteq P}P
    \]
    that is the intersection of all prime ideals that contain the ideal $I$.
\end{proposition}
\begin{proof}
    If $I=R$, then rad$I=I$, and there is nothing to prove. Hence, assume that $I$ is proper, that is $I\neq R$. Then rad$I\subset\bigcap_{P\subset P}P$ (because prime ideals are equal to there radicals (next homework)).
\end{proof}

\chapter{Lecture 17}

\begin{proposition}{}{}
    Let $R$ be a ring, and $S\subset R$. Let $I\subseteq R$ be an ideal, disjoint from $S$. Let $\Sigma_S(I)$ be the set of ideals of $R$ containing $I$, and disjoint from $S$. Then $\Sigma_S(I)$ contains maximal elements. i.e., there exists $J\in\Sigma_S(I)$ such that there is no $J'\in\Sigma_S(I)$ for which $J\subsetneqq J'$.
\end{proposition}
\begin{proof}
    Note that $I\in\Sigma_S(I)$, so $\Sigma_S(I)\neq\phi$. We apply \underline{Zorn's lemma} as follows. Let $b_1\subset b_2\subset\cdots$ be an ascending chain of ideals in $\Sigma_S(I)$ and put
    \[
        b=\bigcup_{j=1}^\infty b_j.
    \]
    $b\in\Sigma_S(I)$ since otherwise there must exists some $s\in S$ such that $s\in b$. But then $s\in b_n$ for some $n\geq 1$ which is a contradiction. Hence $b$ is an upper bound for every chain. Hence, every ascending chain in $\Sigma_S(I)$, so Zorn implies that $\Sigma_S(I)$ has a maximal element. 
\end{proof}
So what is Zorn's lemma?
\begin{lemma}{Zorn's Lemma}{}
    Let $(P,\leq)$ be a partially-ordered set (poset). Suppose that $P\neq\phi$ and every ascending chain in $P$ has an uper bound. Then, $P$ has a maximal element. 
\end{lemma}
\begin{proposition}
    Same hypothesis as the last proposition but also assume that $S$ is multiplicative and $R$ is commutative. Then maximal elements of $\Sigma_S(I)$ are prime ideals. 
\end{proposition}
\begin{proof}
    Let $c\in\Sigma_S(I)$ be maximal. Suppose that $a_1,a_2\in R$ such that $a_1a_2\in C$. If $a_1\in C$ then we are done. Otherwise assume that $a_1\not\in C$. Then $C\subsetneqq C+(a_1)=J_1$. And since $C$ is maximal then $J_1\not\in\Sigma_S(I)$. Hence, $S\cap J_1\neq\phi$, say $s_1\in S\cap J_1$. $s\in J$ implies $s_1=c_1+a_1r_1$ for $c_1\in C,r_1\in R$. Similarily, we must assume $a_2\not\in C$. Hence, $J_2=C+(a_2)$ properly contains $C$, so $S\cup J_2\neq\phi$, say $s_in S\cup J_2$. Then $s_2=c_2+a_2r_2$, for $c_2\in C$ and $r_2\in R$. Then 
    \[
        \underbrace{s_1s_2}_{\in S}=(c_1+a_1r_1)(c_2+a_2r_2)=c_1c_2+c_1a_2r_2+c_2a_1r_1+a_1r_1a_2r_2\in C.
    \]
    This is a contradiction sicne we assumed that $S$ and $C$ are disjoint. Thus, maximal ideals are prime ideals. 
\end{proof}
This proposition implies  
\begin{proposition}{}{}
    \[
        \text{rad(rad($I$))}=\text{rad($I$)}.
    \]
\end{proposition}
\begin{proof}
    $S=\{1,a,a^2,\dots\}$, $S$ is multiplicative and $a\not\in$rad$I$. The proposition implies there is a prime ideal $P$ containing $I$ and disjoint from $S$. Hence, $a\not\in\cap_{I\subset P,\text{ prime}}P$.
\end{proof}
\begin{definition}{Jacobson Radical}{}
    Let $R$ be a ring. The Jacobson Radical of $R$ is the ideal 
    \[
        \mathcal{J}(R)=\bigcap_{\substack{m\text{ maximal}\\\text{ideal of $R$}}}m.
    \]    
\end{definition}
\begin{proposition}
        Let $R$ be a unital ring. Then $a\in\mathcal{J}(R)$ \textit{iff} $1-ac$ is a unit for all $c\in R$. 
\end{proposition}
\begin{proof}
    It suffices to show that there exists a maximal ideal $m$ not containing $a$ \textit{iff} there exists $c\in R$ such that $1-ac$ is not a unit. \\
    First suppose that $a\not\in m$, $m$ maximal. Then $m+(a)=R$. Then $1=m_1+ab$, with $m_1\in M$ and $b\in R$. Hence, $m_1=1-ab$, and $m_1$ is not a unit. Conversely, suppose that $1-ac$ is not a unit. Then $(1-ac)$ is a proper ideal of $R$. Thus, $(1-ac)\subset m$ for some maximal ideal $m$. We claim that $a\not\in m$. This si because if otherwise then $m_1=1-ac\implies1\in m$, which is a contradiction. 
\end{proof}
\addunumberedsubsec{Pell's Equation}
\[
    x^2-2y^2=1
\]
with $x,y\in\Z_{\geq1}$. Euler wanted to solve 
\[
    \frac{m(m+1)}{2}=n^2
\]
with the LHS being triangular numbers. He got
\begin{align*}
    4m(m+1)&=8n^2\\
    4m^2+4m+1&=8n^2\\
    (2m+1)^2-2(2n)^2&=1.
\end{align*}
Solving Pell's equation is the same as finding all units in $\Z[\sqrt{2}]$.

\chapter{Lecture 18}
\begin{proposition}{Prime Avoidance}{}    
    Let $\mathcal{P}_1,\dots,\mathcal{P}_r$ be prime ideals in a ring $R$. If $I\subsetneq \mathcal{P}_i$ ($I$ an ideal of $R$) for $i=1,\dots,r$, then $I\subsetneqq\bigcup_{j=1}^r\mathcal{P}_i$.
\end{proposition}
\begin{proof}
    If $r=1$, the claim is obvious. Hence, we assume that $r\geq2.$ Suppose that $I\subset I\subsetneqq\bigcup_{j=1}^r\mathcal{P}_i$, but $I$ is not contained in the union of fewer $\mathcal{P}_i$'s. In particular for $1\leq i\leq r$ $I\subsetneqq\bigcup_{j\neq i}\mathcal{P}_j$. Hence there exists $a_i\in I\setminus\bigcup_{j\neq i}\mathcal{P}_j$. But then $a_i\in\mathcal{P}_i$. Take $a=a_1\cdots a_{r-1}+a_r\in I$.\\
    \textit{Claim:} $a\not\in\mathcal{P}_i$ for $1\leq i\leq r$. To see this note that $a_1,\dots,a_{r-1}\not\in\mathcal{P}_r$.\\
    Since $\mathcal{P}_r$ is prime, their product must not lie in $\mathcal{P}_r$. By definition, $a_r\in\mathcal{P}_r$. Hence $a\not\in\mathcal{P}_r$. If $1\leq i\leq r-1$, then $a_1\cdots a_{r-1}\in\mathcal{P}_i$, but $a_r\not\in\mathcal{P}_i$. $a\in I$, but $a\not\in\mathcal{P}_i$ for $1\leq i\leq r$, contradicts $I\subset\bigcup_{j=1}^r\mathcal{P}_j$.
\end{proof}
$\Z[\sqrt{-5}]$ is a quadratic extention of the $\Z$ (the ring of integers extended with roots for unsolvable (monoic) quadratic polynomials). $(5)\in\Z$ is irreducible. But $(5)=(\sqrt{-5})\in\Z[\sqrt{-5}]$ is reducible.\\
Let $\phi:A\to B$ be a ring homomorphism ($B$ has the structure of an $A$-algebra). If $b$ is an ideal in $B$, then $\phi^{-1}b$ is an ideal in $A$. 
\begin{lemma}{}{}
    If $b$ is an ideal in $B$, then $\phi^{-1}b$ is an ideal in $A$.\\
\end{lemma}
\begin{proof}
    \textit{Recall:}
    \[
        \phi^{-1}b=\{a\in A:\phi(a)\in b\},
    \]
    $0_A\in\phi^{-1}(\{0_B\})$, so we need to check 
    \begin{enumerate}
        \item $0_B\in b$
        \item If $a_1,a_2\in\phi^{-1}(b)$ then there exists $b_1,b_2\in b$ such that $a_i\in\phi^{-1}(\{b_i\})$, (i.e., $b_i=\phi(\{a_i\})\in b$). $b_1+b_2\in b$, $b$ is an ideal. $b_1+b_2\phi(a_1)+\phi(a_2)=\phi(a_1+a_2)\Rightarrow a_1+a_2\in\phi^{-1}(b)$. Supose $r\in A$, $a\in\phi^{-1}(b)$. Then $\phi(ra)=\phi(r)\phi(a)$. $\phi(a)\in b$, hence $\phi(r)\phi(a)\in b$. Hence $\phi(ra)\in b$ and $ra\in\phi^{-1}(b)$.
    \end{enumerate}
    Thus, we have our result.
\end{proof}
\begin{definition}{Ideal Contractions}{}
    Let $\phi:A\to B$, be a ring homomorphism. For an ideal $b$ of $B$, $\phi^{-1}(b)$ is called the \underline{contraction} of $b$ in $A$.
\end{definition}
\begin{example}{}{}
    $\Z[\sqrt{-5}]$ contracts to what in $\Z$?\\
    $\phi^{-1}(\Z[\sqrt{-5}])=(6)$.
\end{example}
In general $\phi(I)$ is \underline{not} an ideal in $B$. But we can take the ideal generated by $\phi(I)$ and call that the \underline{extension} of $I$.
\begin{lemma}{}{}
    If $A$ is a subring of $B$, and $b\subset B$ an ideal ($b^c$ is the contraction of $b$ in $A$), then $b^c=b\cap A$ ($\phi:A\to A$ is the identity map).
\end{lemma}
\begin{proof}
    $a\in b^c=\phi^{-1}(b)$ \textit{iff} $\phi(a)\in b$, but $\phi(a)=a$, so $a\in b$. Hence $b^c=b\cap A$.
\end{proof}
\addunumberedsubsec{Chinese Remainder Theorem}
\textit{Classical:} If $m_1,\dots,m_k$ are pairwise coprime integers. and $a_1,\dots,a_k$ are arbitrary integers, then there exists $x\in\Z_{\geq1}$ such that $x-a_i$ is divisible by $m_i$ for $i=1,\dots,k$.\\
We have to interpret the second half of the statement as a statement about the map $\phi:R\to R/m_1\times\cdots R/m_k$, $a\mapsto(a+m_1,\dots,a+m_k)$.

\chapter{Lecture 19}
\begin{theorem}{Chinese Remainder Theorem for Rings}{}
    Let $R$ be a unital ring, and $m_1,\dots,m_k$ are pairwise co-prime ideals in $R$. Then the map
    \[
        \phi:R\to R/m_1\times R/m_2\times\cdots\times R/m_k
    \]
    where
    \[
        \phi(a)=(a+m_1,\dots,a+m_k)
    \]
    is a surjective ring homomorphism, and 
    \[
        \ker(\phi)=(m_1\cdot m_2\cdots m_k)
    \]
\end{theorem}
\begin{proof}
    By induction on $k$.\\
    $k=2$: Since $(m_1+m_2)=R$, there exists $a_i\in m_i$ such that $a_1+a_2=1$. Then for any $x_1,x_2\in R$ we have $\phi(a_1x_2+a_2x_1)=(x_1+m_1,x_2+m_2)$. This shows that $\phi$ is surjective.
    \begin{lemma}{}{}
        For $I,J$ prime ideals in $R$, $I\cdot J=I\cap J$
    \end{lemma}
    \begin{proof}
        $I\cdot J=\{i_1j_1+\cdots i_kj_k:i_\ell\in I,j_\ell\in J\}$. Since $I,J$ are coprime, there exists $i\in I$ and $j\in J$ such that $i+j=1$. If $a\in I\cap J$, then $a(i+j)=a\in I\cdot J$. Hence, $I\cap J\subseteq I\cdot J$.\\
        Conversely, suppose that $a=i_1j_1+\cdots i_kj_k\in I\cdot J$. Then each $i_\ell j_\ell\in I\cap J$, so $a\in I\cap J$. Hence, $I\cdot J\subseteq I\cap J$. Giving the desired result. 
    \end{proof}
    Suppose $c\in m_1\cap m_2=m_1\cdot m_2$. Then $c(a_1+a_2)=c\in m_1\cdot m_2$ and
    \[
        \phi(c)=\phi(a_1c+a_2c)=(c+m_1,c+m_2)=(0,0).
    \]
    Hence $c\in\ker(\phi)$. If $c\in\ker(\phi)$, then $\phi(c)=(0+m_1,0+m_2)$. But $c=a_1c+a_2c\Rightarrow\phi(c)=(c+m_1,c+m_2)$. Hence $c\in m_1$ and $c\in m_2$. That is, $c\in m_1\cap m_2=m_1\cdot m_2$.\\
    \textit{Inductive Hypothesis:} Let $k\geq2$. Suppose that the statement holds for $k$. We need to show that it is true for $k+1$.\\
    We have $m_1,m_2,\dots,m_k,m_{k+1}$ pairwise co-prime ideals. Take $M=m_1\cdots m_k$. Have to show that $M+m_{k+1}=R$. By the hyopthesis, there exist $a_i\in m_{k+1}$ and $b_i\in m_i$ such that $a_1+b_i=1$. Let
    \[
        \alpha=\prod_{1\leq i\leq k}\underbrace{(a_i+b_i)}_{=1}. 
    \]
    Then $\alpha\in m_{k+1}+m_1+\cdots+m_k$. Hence, $a\in m_{k+1}+M$, so $R=m_{k+1}+M$. Hence, using the $k=2$ case
    \[
        \phi:R\to R/M+R/m_{k+1}
    \]
    is surjective. By the induction hypothesis, $R/M\cong R/m_1+\cdots R/m_k$, so we have our desired result. 
\end{proof}
\addunumberedsubsec{Noetherian Rings}
Recall that a ring $R$ is Noetherian is it satifies the ascending chain condition. 
\begin{proposition}{}{}
    Let $R$ be a ring. The following are equivalent for an $R$-module $M$
    \begin{enumerate}[label=(\alph*)]
        \item Every submomdule of $M$ is finitely generated;
        \item Every ascending chain of submodules $M_1\subset M_2\subset\cdots$ in $M$ is eventually constant;
        \item Every non-empty set of submodules of $M$ has a maximal element. 
    \end{enumerate}
\end{proposition}
\begin{proof}
    ((a)$\Rightarrow$(b)) Suppose $M_1\subset\cdots\subset M_k\subset\cdots$ is an ascending chain of modules. Then, just like the ideal case, one shows that $\bigcup_{j\geq1}M_j=N$ is a submodule of $M$. By (a), $N$ is finitely generated, say $N=\langle n_1,\dots,n_\ell\rangle$. There exists $L\geq1$ such that $n_j\in M_j$ for $1\leq j\leq \ell$. Hence, $M_t\subseteq M_L$ for all $t\geq1$, so $M_t=M_L$ for $t\geq L$.\\
    ((b)$\Rightarrow$(c)) Let $\Sigma$ be a non-empty set of submodules. If $\Sigma$ has no maximal element, then by the contrapositive of Zorn's lemma, there must be ascending chains with no maximal element in $\sigma$, contradiciting the hypothesis.\\
    ((c)$\Rightarrow$(a)) Let $N$ be a submodule of $M$ and let $\Sigma(N)$ be the set of finitely generated submodules contained in $N$. $\Sigma(N)$ is non-empty (since it contains $\{0\}$), hence $\Sigma(N)$ contains a maximal element, say $N'=\langle n_1,\dots,n_\ell\rangle$. If $N'\neq N$, then there exists $n\in N\setminus N'$. Then $N''=\langle n_1,\dots,n_\ell,n\rangle$ properly contains $N$. But $N''\subset N$, hence $N''\in\Sigma(N)$. Thus $N''$ is larger than $N'$ contradiciting the maximality of $N'$.
\end{proof}

\chapter{Lecture 20}
\addunumberedsubsec{Exact Sequences}
Exact sequences are of the form
\[
    0\xrightarrow{f_0}A_1\xrightarrow{f_1}\cdots\xrightarrow{f_{k-1}}A_k\xrightarrow{f_k}0
\]
such that Im$(f_j)=\ker(f_{j+1})$
\addunumberedsubsec{Short Exact Sequence}
\[ 
    0\xrightarrow{f_0}A\xrightarrow{f_1}B\xrightarrow{f_2}C\xrightarrow{f_3}0
\]
\begin{example}{}{}
    \[
        0\xrightarrow{f_0}\R\xrightarrow{f_1} C^\infty(\R,\R)\xrightarrow{f_2=\frac{d}{dx}} C^\infty(\R,\R)\to 0
    \]
    $f_1:\R\mapsto$ constant functions. 
\end{example}
\begin{example}{}{}
    \[
        0\to\Z\to\Q\to\Q/\Z\to 0
    \]
\end{example}
\begin{example}{}{}
    \[
        0\to\Z\to\R\to\R/\Z\cong S^1(\text{sphere})\to0
    \]
\end{example}
\begin{example}{}{}
    Selmer group
    \[
        0\to E^{(\Q)}/2E(\Q)\to\text{Sel}_2(E)\to\text{Sha}_E[2]\to0
    \]
    $E$ is an elliptic curve over $\Q$. Computing Sha is an open problem for elliptic curves.
\end{example}
\addunumberedsubsec{Commutative Diagrams}
\[
    \begin{tikzcd}[column sep=huge, row sep=huge]
        \mathcal{A} \arrow[r, "F"] \arrow[d, "\alpha"'] & \mathcal{B} \\
        \mathcal{C} \arrow[r, "G"'] & \mathcal{D} \arrow[u, "\beta"']
    \end{tikzcd}
\]
with $F=\beta\circ G\circ\alpha$
\begin{example}{}{}
    \[
        \begin{tikzcd}[column sep=large, row sep=large]
            % Row 1: A, B, (empty), (empty)
            X(\Q) \arrow[r, ] \arrow[d, ] & x(\Q_p) \arrow[d, ] & & \\
            % Row 2: C, D, (empty), E
            J(\Q) \arrow[r, ] & J(\Q_p) \arrow[rr, ] & & H^\circ(x_{\Q_p}\Omega')
            \arrow[from=1-2, to=2-4, , bend left=15]
        \end{tikzcd}
    \]
\end{example}
\begin{proposition}{}{}
    Let $R$ be a ring, and let 
    \[ 
        0\xrightarrow{f_0}M_1\xrightarrow{f_1}M_2\xrightarrow{f_2}M_3\xrightarrow{f_3}0
    \]
    be an exact sequence of $R$ modules. 
    \begin{enumerate}[label=(\alph*)]
        \item If $N\subseteq P$ are submodules of $M_2$ such that $f_1(M_1)\cap N=f_1(M_1)\cap P$ and $f_2(N)=f_2(P)$, then $N=P$;
        \item If $M_1$ and $M_1$ are finitely generaeted, then so it $M_2$;
        \item $M_2$ is Noetherian \textit{iff} $M_1$ and $M_3$ are both Noetherian.
    \end{enumerate}
\end{proposition}
\begin{proof}
    \begin{enumerate}[label=(\alph*)]
        \item Let $p\in P$. The second condition implies there exists $n\in N$ such that $f_2(p)=f_2(n)$. Since $f_2$ is a module homomorphism we have $f(n-p)=0$. Therefore, sicne this si an exact sequence, $n-p\in\ker(f_2)\Rightarrow n-p\in$Im$f_1=(f_1(M_1))$. Since $N\subseteq P$ we have $n,\in P$, hence $n-p\in P$. Therefore, $n-p\in f_1(M_1)\cap P$. By the hypothesis $f_1(M_1)\cap P=f_1(M_1)\cap N$. Hence, $n-p\in f_1(M_1)\cap N$, thus $s-n\in N$. But $\underbrace{(p-n)}_{\in N}+\underbrace{n}_{\in N}=p\in N$. Thus, $P\subseteq N$. So $P=N$.
        \item Let $s_1$ be a finite set of generators of $M_1$. Recall, from the property of exact sequences that $f_2:M_2\to M_3$ is surrjective. Since $M_3$ is finitely generated, there exists finite $s_2\subset M_2$ such that $f_2(s_3)$ generates $M_3$. Let $N\subset M_2$ be the submodules generated by $f_1(s_1)$ and $s_3$. Then $f_1(M_1)\cap N=M_2$, and $f_2(N)=M_3=f_2(M_2)$. Part (a) says that $N=M_2$. Hence $M_2$ is generated by $f_1(s_1)\cup s_3$, which is finite.
        \item One direction it clear (i.e., $M_1,M_2$ are Noetherian implies that $M_2$ is Noetherian). Now assume $M_2$ is Noetherian. Let
        \[
            N_1^{(1)}\subseteq N_2^{(1)}\subseteq\cdots
        \]
        be an ascending chain of submodiles of $M_1$. then
        \[
            f_1(N_1^{(1)})\subseteq f_1(N_2^{(1)})\subseteq\cdots 
        \]
        is an ascending chain in $M_2$. Since $M_2$ is Noetherian there exists $k\leq 1$ such that $f_1(N_s^{(1)})=f_1(N_k^{(1)})$ for all $s\geq k$. But $f_1$ is injective, so this implies 
        \[
            N_s^{(1)}=N_k^{(1)}
        \]
        for all $s\geq k$. So $M_1$ is Noetherian.\qedhere
    \end{enumerate}
\end{proof}

\chapter{Lecture 21}
\addsec{Grothendieck's proof of Riemann-Roch}
\begin{enumerate}
    \item Introduced the notion of \underline{schemes} rather than varieties;
    \item Introduced the idea that Riemann-Roch is about \underline{morphisms}, not the actual varieties;
    \item Decomposing morphisms in a certain canonical manner gives ``obvious'' proof.
\end{enumerate}
\begin{proposition}{}{}
    Every finitely generated module over a Noetherian ring is Noetherian.
\end{proposition}
\begin{proof}
    Let $M$ be a finitely generated module over a Noetherian ring $R$. If $M$ is generated by a single elements over $R$, say $\alpha$, then $M\cong R/I$ for some ideal $I\subset R$. Since $R$ is Noetherian, so is $R/I$, so we're done. Now induct on teh number $n$ of generators of $M$, say $M=\langle a_1,\dots,a_n \rangle$. Let $N=\langle a_1,\dots,a_{n-1} \rangle$ by the submodule generated by $a_1,\dots,a_{n-1}$. By the IH, $N$ is Noetherian. THen $M/N$ is generated by one element; base case finishes the job. 
\end{proof}
\begin{proposition}{}{}
    Every finitely generated module $M$ over a Noetherian ring $R$ contains a finite chain of submodules $M_r\subset\cdots\subset M_1\subset M_0$ such that $M_i/M_{i+1}$ is isomorphic to $R/P_i$ for some prime ideal $P_i$ or $R$.
\end{proposition}
The proof needs the idea of annihilator.
\begin{definition}{Annihilator}{}
    Let $M$ be an $R$-module. For $x\in M$ the \underline{annihilator} of $x$, denoted ann$(x)$, is the set ann$(x)=\{r\in R:rx=0\}$. 
\end{definition}
\begin{lemma}{}{}
    ann($x$) is an ideal in $R$ ($R$ is commutative).
\end{lemma}
\begin{proof}
    If $r_1,r_2\in\ann(x)$, then $(r_1+r_2)x=r_1x+r_2x=0$. If $r_1\in\ann(x)$, $r_2\in R$ then $(r_2r_1)x=r_20=0=(r_1r_2)x$.
\end{proof}
\begin{proof}[Proof of Prop 1.2]
    Let $\mathfrak{a}=\ann(x)$ be maximal among the annihilators of non-zero elements of $M$. We claim that $\ann(x)$ is prime. Suppose $a_1a_2\in\mathfrak{a}$. Then $(a_1a_2)x=0$. Then 
    \[
        \mathfrak{a}\subset(a_1)+\mathfrak{a}\subset \ann(a_2x).
    \]
    If $a_2\in\mathfrak{a}$, then $a_2x\neq0$. Hence, $\mathfrak{a}=\ann(a_2x)$ by maximality. This implies $a_1\in\mathfrak{a}$. To procees, nore that the submodule $Rx$ (orbit of $x$ acted on by $R$) is isomorphic to $R/\ann(x)$. If $M$ is non-zero, then there exists a non-zero $x\in M$ sucj that $\ann(x)$ is maximal, hence prime. Hence $M$ contains a submodule $M_1=Rx$ which is isomorphic to $R/\ann(x)=R/P_1$. Similarily, $M_2$ contains a submodule isomorphic to $R/P_2$. This terminates because $M$ is Noetherian. 
\end{proof}
\begin{theorem}{Hilbert Basis Theorem}{}
    Every finitely generated algebra over a Noetherian ring $R$ is Noetherian. 
\end{theorem}
\begin{proof}[Proof of the Theorem]
    Let $A$ be a ring, $B$ an $A$-alegbra. We argue by induction over minimum number of generators (for $B$). Note $A[x_1,\dots,x_n]\cong A[x_1,\dots,x_{n-1}][x_n]$. Thus, only the case $n=1$ matters. Let $f(x)=c_rx^r+\cdots+c_1x+c_0\in A[x]$, and $I$ be an ideal of $A{x}$. Let $I(k)$ be the set of elements in $A$ which appear as the leading coefficients of some $f\in I$ with $\deg(f)\leq k$ (can be 0). Then $I(k)$ is an ideal in $A$. Further, $I(k-1)\subset I(k)$, since $c_{k-1}x^{k-1}+\cdots+c_1x+c_0\in I$ implies that $c_{k-1}x^{k}+\cdots+c_1x^2+c_0x\in I$. If $J$ is an ideal in $A[x]$ containing $I$, then $I(k)\subset J(k)$. Further if $I(k)=J(k)$ for all $k\geq0$ then $I=J$. Hence, suppose there exists $f\in J\setminus I$, and $f$ has a minimal degree among all such polynomials. Let $k$ be this degree. Then $I(k)=J(k)$, because for all $g\in J$ of degree at most $k-1$, $g\in I$. If one takes such a $g$ then $f+g\in I$. Then the leading coefficient of $f$, and $f+g$ are the same. Continue on like this\\
    Since $A$ is Noetherian
    \[
        I(1)\subset I(2)\subset \cdots
    \]
    must terminate, say at $I(k)=I(\ell)$ for all $k\geq\ell$. For $i\leq\ell$, there exists a finite generating set $S_i=\{C_{i1},C_{i2},\dots\}$ for $I(i)$. For each $(i,j)$ there exists $f_{ij}\in I$ for degree $i$ and leading coefficient be $C_{ij}$. Let $J=\langle C_{ij} \rangle$. Further, $J(k)=I(k)$ for every $k$, so $I=J$. 
\end{proof}
\chapter{Lecture 22}
\begin{lemma}{Nakayama's Lemma}{}
    Let $R$ be a ring. Let $I\subset R$ be an ideal, $M$ a finitely generated $R$-module. Suppose $I\subset m$ for all maximal ideals $m$ of $R$. Then:
    \begin{enumerate}[label=(\alph*)]
        \item If $M=I\cdot M$, then $M=\{0\}$;
        \item If $N$ is a submodule of $M$ such that $M=N+I\cdot M$, then $M=N$.
    \end{enumerate}
\end{lemma}
Nakayama's lemma is most applicable to \underline{local rings}, which are rings with a unique maximal ideal. Rings with non-trivial Jacobson ideals. 
\begin{proof}[Proof of Lemma 1.1]
    (a) Suppose $M\neq\{0\}$. Choose a minimal set of generators $\{a_1,\dots,a_k\}$. Since $M=I\cdot M$, we have a relation of the form $a_1=u_1a_1+\cdots+u_ka_k$, $u_j\in I$ for $1\leq j\leq k$. Then $(1-u_1)a_1=u_2a_2+\cdots+u_ka_k$. We claim that $1-u_1$ cannot lie in any maximal ideal $m$. Suppose otherwise, then $1-u_1=v\in m$, so $1=u_1+v\in m$ (since $I\subset m$). This contradicts $m$ being a proper ideal. Hence, $1-u_1$, must be a unit ($u_1\neq0$). In particular $a_1\in\langle a_2,\dots,a_k \rangle$. This contradicts the minimality of $\{a_1,\dots,a_k\}$ as a generating set. \\
    (b) If $M=N+I\cdot M$, then $M/N\cong I(M/N)$. By part (a) we have $M/N=\{0\}$.
\end{proof}
\begin{corollary}{}{}
    Let $A$ be a local ring with a maximal ideal $m$. Let $F=A/m$ be the \underline{residue field} of $A$. Let $M$ be a fintely generated module over $A$. Then action of $A$ on $M/mM$ \underline{factors through} $F$, and elements $a_1,\dots,a_k$ in $M$ generate $M$ as an $A$-module \textit{iff} $a_i+mM$, $1\leq i\leq k$, generate $M/m$ as a \underline{$F$-vector-space}. Factors through means we can break the action into two parts which go through $F$.
\end{corollary}
\begin{proof}
   If $M=\langle a_1,\dots,a_k \rangle$, then it is clear that $M/m=\langle a_1+mM,\dots,a_k+mM \rangle$.\\
   Now suppose $M/m=\spn_F(a_1+mM,\dots,a_k+mM)$. Let $N=\langle a_1,\dots,a_k \rangle$. Then 
   \[
    N\underbrace{\to}_{\text{embedding}} M\underbrace{\to}_{\text{surjective}} M/mM.
   \]
   Hence, $M=N+mM$. Nakayama $\Rightarrow$ $M=N$.
\end{proof}
\begin{corollary}
    Let $A$ be a Noetherian local ring with a maximal ideal $m$. Elements $a_1,\dots,a_k$ generate $m$ as an ideal \textit{iff} $a_1+m^1,\dots,a_k+m^2$ generate $m/m^2$ as a vector space over $F=A/m$. In particular, the minimum number of generators of $m$ as an ideal equal $\dim_F(m/m^2)$.
\end{corollary}
\begin{proof}
    Because $A$ is Noetherian, $m$ is finitely generated. Any ideal of $A$ is an $A$-module, apply the previoud result with $M=m$. 
\end{proof}
\addunumberedsubsec{Dedekind Domains}
Integral closure of a domain.\\
Let $A$ be a unital, commutative domain (no zero divisors). Suppose $A\subset B$, $B$ is also a ring (that is $A$ is a subring of $B$). We say that $b\in B$ is \underline{integral} over $A$ if there exists a \underline{monic} polynomial in $A[x]$ say $f(x)=x^r+c_{r-1}x^{r-1}+\cdots+c_1x+c_0$ such that $f(b)=0$. The integral closure $\overline{A}$ of $A$ in $B$ is the set of all elements in $B$ which are integral over $A$. (Note that $A\subset \overline{A}$). 
\begin{definition}{Dedekind Domain}{}
    An integral domain not equal to a field is said to be a Dedekind domain if 
    \begin{enumerate}[label=(\alph*)]
        \item $A$ is Noetherian;
        \item $A$ is equal to its integral closure in the field of fractions;
        \item Every non-zero prime ideal of $A$ is maximal. 
    \end{enumerate}
\end{definition}
\begin{theorem}{}{}
    Let $\alpha$ be an algebraic integer (is a zero of a polynomial in $\Z$). Then the ring $\Z[\alpha]$ is a Dedekind domain. 
\end{theorem}
\begin{theorem}{Dedekind}{}
    Let $A$ be a Dedekind domain domain, and $I\subset A$ an ideal. Then $I$ admits a \underline{unique factorization}
    \[
        I=P_1^{a_1}\cdots P_K^{a_k}
    \]
    into distinct prime ideals. The $P_1,\dots,P_k$ are precisely the prime (max) ideals containing $I$.
\end{theorem}
The ideals $A$ and $B$ of a ring $R$ are coprime if $A+B=R$.
\end{document}